\documentclass[11pt, a4paper]{amsart}
\usepackage{amsmath, amssymb, amsthm, amsfonts}
\usepackage{mathtools}
\usepackage[margin=1in]{geometry}
\usepackage{hyperref}
\usepackage{enumitem}

\newtheorem{theorem}{Theorem}[section]
\newtheorem{lemma}[theorem]{Lemma}
\newtheorem{proposition}[theorem]{Proposition}
\newtheorem{corollary}[theorem]{Corollary}
\newtheorem{conjecture}[theorem]{Conjecture}
\theoremstyle{definition}
\newtheorem{definition}[theorem]{Definition}
\newtheorem{remark}[theorem]{Remark}
\newtheorem{hypothesis}[theorem]{Hypothesis}
\newtheorem{example}[theorem]{Example}
\newtheorem{notation}[theorem]{Notation}

\newcommand{\R}{\mathbb{R}}
\newcommand{\N}{\mathbb{N}}
\newcommand{\Q}{\mathbb{Q}}
\newcommand{\C}{\mathbb{C}}
\newcommand{\Z}{\mathbb{Z}}
\newcommand{\OO}{\mathbb{O}}
\newcommand{\HH}{\mathbb{H}}
\newcommand{\pphi}{\varphi}
\newcommand{\SG}{\Sigma}
\newcommand{\cas}{{\rm cas}}
\newcommand{\met}{{\rm met}}
\newcommand{\hyd}{{\rm hyd}}
\newcommand{\gge}{{\rm gauge}}
\newcommand{\co}{{\rm coh}}
\newcommand{\TM}{{\rm TM}}
\newcommand{\CTM}{{\rm CTM}}
\newcommand{\Fix}{\operatorname{Fix}}
\newcommand{\Der}{\operatorname{Der}}
\newcommand{\Met}{\operatorname{Met}}
\newcommand{\Conn}{\operatorname{Conn}}
\newcommand{\Hol}{\operatorname{Hol}}
\newcommand{\Sym}{\operatorname{Sym}}
\newcommand{\Sp}{\operatorname{Sp}}

\title{Cascade--Classical Correspondence: Foundations}
\author{}
\date{}

\begin{document}

\maketitle

\begin{abstract}
This paper constructs the \emph{correspondence layer} between the
cascade framework \cite{CascadeFoundations} and classical
mathematical structures. Seven sectorwise projections are defined
as precise mathematical maps with specified domains, codomains,
measurability, continuity, and symmetry-equivariance: the metric
projection $\Pi_{\met}$, the hydrodynamic projection $\Pi_{\hyd}$,
the gauge projection $\Pi_{\gge}$, the regularized spectral-zeta
transfer operator $T_\zeta[h]$, the regularized Hecke transfer
operator $T_E[h]$, the Galois fixed-part projector $\Pi_\sigma$,
and the cascade-to-Turing transport pair $(\Pi_{\TM}, I_{\TM})$.
For each projection we establish either an unconditional
well-definedness and equivariance theorem, or a conditional
correspondence theorem under one of four explicitly named open
hypothesis families
$\mathbf{H_1}$ (geometric limit),
$\mathbf{H_2}$ (gauge reconstruction),
$\mathbf{H_3}$ (spectral--arithmetic trace), and
$\mathbf{H_4}$ (universality / normal form).

The paper does \emph{not} prove any of the six open Clay Millennium
Prize Problems. Its purpose is narrower and structural: it
identifies the informal ``cascade-to-classical projection'' language
used across the six companion manuscripts and replaces it with
precise, Clay-bar-verifiable mathematics, so that the companion
papers for Yang--Mills, the Riemann Hypothesis, Navier--Stokes, the
Hodge Conjecture, the Birch--Swinnerton-Dyer Conjecture, and $P$
vs $NP$ can be formulated against a stable foundation. Each
companion paper inherits exactly one of
$\mathbf{H_1}$--$\mathbf{H_4}$; what remains open after that
inheritance is, in each case, a clearly-stated classical analytic
or algorithmic problem, not a cascade-internal ambiguity.

We audit the upstream cascade literature explicitly in
Section~\ref{sec:audit}, classifying each cited result as
\emph{theorem-grade}, \emph{formal framework}, or
\emph{motivational only}. Load-bearing claims in this paper rely
only on theorem-grade cascade inputs or on classical mathematics.
Items of lower upstream status are either reproved here from
explicit formulas or relegated to open-hypothesis status.
\end{abstract}

\tableofcontents

%=====================================================================
\section{$\mathbf{F1}$: The Permeability Fixed-Point Axiom}
\label{sec:F1}
%=====================================================================

\noindent\emph{This section states, at Clay-bar standard, the single
non-classical axiom of the cascade-correspondence programme.
Downstream papers in the $\alpha$-chain
(\texttt{cascade-12d-closure.tex},
\texttt{cascade-phason-coxeter.tex},
\texttt{cascade-fine-structure.tex}) cite this section as
``\S F1'' and import the axiom \textbf{only}, not any subsequent
machinery of this manuscript. No argument in \S F1 depends on
any other section of this paper.}

\begin{notation}[Permeability parameter]
\label{not:r}
Throughout \S F1, $r$ denotes a positive real variable.
\end{notation}

\begin{axiom}[$\mathbf{F1}$: permeability fixed-point axiom]
\label{ax:F1}
There exists a positive real $r \in (0, \infty)$ satisfying
\begin{equation}
\label{eq:F1}
  r \;=\; 1 \;+\; \frac{1}{r}.
\end{equation}
This equation, and the existence of a positive solution, is
the sole non-classical axiom of the cascade-correspondence
programme. We refer to any such positive solution as a
\emph{permeability fixed point}.
\end{axiom}

\begin{proposition}[Uniqueness and explicit value]
\label{prop:F1-explicit}
There exists a unique positive real $r \in (0, \infty)$
satisfying $\mathbf{F1}$ of Axiom~\ref{ax:F1}, and it equals
the golden ratio
\[
  r \;=\; \varphi \;:=\; \frac{1 + \sqrt{5}}{2}.
\]
\end{proposition}

\begin{proof}
Rearranging equation~\eqref{eq:F1}:
\[
  r \;=\; 1 + \frac{1}{r}
  \;\iff\;
  r^2 \;=\; r + 1
  \;\iff\;
  r^2 - r - 1 \;=\; 0.
\]
The quadratic $r^2 - r - 1$ has discriminant $1 + 4 = 5$
and roots $(1 \pm \sqrt{5})/2$. Only $(1 + \sqrt{5})/2$ is
positive (since $\sqrt{5} > 1$). Hence the unique positive
real satisfying $\mathbf{F1}$ is $r = (1 + \sqrt{5})/2 =
\varphi$.
\end{proof}

\begin{corollary}[Quadratic and self-similarity identities]
\label{cor:F1-identities}
The golden ratio $\varphi = (1 + \sqrt{5})/2$ satisfies:
\begin{align}
  \varphi^2 &\;=\; \varphi + 1, \label{eq:phi-sq} \\
  \varphi^{-1} &\;=\; \varphi - 1, \label{eq:phi-inv} \\
  \varphi &\;=\; 1 + \frac{1}{\varphi}. \label{eq:phi-recursive}
\end{align}
\end{corollary}

\begin{proof}
Equation~\eqref{eq:phi-sq} is an immediate restatement of
$r^2 - r - 1 = 0$ from the proof of
Proposition~\ref{prop:F1-explicit}. Dividing $\varphi^2 =
\varphi + 1$ by $\varphi$ (which is nonzero) gives $\varphi
= 1 + 1/\varphi$, establishing~\eqref{eq:phi-recursive};
equivalently, $\varphi - 1 = 1/\varphi$, establishing
\eqref{eq:phi-inv}.
\end{proof}

\begin{remark}[Brouwer existence and eventual contraction]
\label{rmk:F1-banach}
The equation \eqref{eq:F1} can be read dynamically as a
fixed-point problem for the map
$f \colon (0, \infty) \to (0, \infty)$, $f(r) := 1 + 1/r$.
One verifies directly that $f$ maps $[1, 2]$ into $[3/2, 2]
\subset [1, 2]$ (since $r \in [1, 2]$ implies $1/r \in
[1/2, 1]$, hence $f(r) \in [3/2, 2]$), so $f$ restricts to a
continuous self-map of the closed interval $[1, 2]$ and
Brouwer's fixed-point theorem provides existence of a fixed
point.

Iterative convergence follows by an eventual-contraction
argument: on the one-step image $[3/2, 2]$ the derivative
bound $|f'(r)| = 1/r^2 \leq 4/9 < 1$ holds uniformly, so the
restriction $f|_{[3/2, 2]} \to [3/2, 2]$ is a $(4/9)$-Lipschitz
contraction; the Banach fixed-point theorem then applies from
the first iterate onward. Any iteration $r_{n+1} = f(r_n)$
starting from $r_0 \in [1, 2]$ has $r_1 \in [3/2, 2]$ and
thereafter converges to the unique fixed point
Proposition~\ref{prop:F1-explicit} identifies.

For the purposes of \S F1, the existence claim in
Axiom~\ref{ax:F1} is delivered independently and explicitly
by Proposition~\ref{prop:F1-explicit} (via the quadratic
formula), so no Brouwer / Banach machinery is load-bearing
for downstream papers.
\end{remark}

\begin{remark}[Scope of $\mathbf{F1}$]
\label{rmk:F1-scope}
\S F1 makes \emph{no} claim beyond the existence/uniqueness
of the positive solution $\varphi$ to $r = 1 + 1/r$. In
particular, \S F1 does not assert:
\begin{itemize}
\item any physical interpretation of ``permeability'';
\item any connection to cascade sectors, $E_8$, $H_4$, the
  $600$-cell, or Coxeter theory;
\item any connection to the fine-structure constant, the
  Standard Model, or any Millennium problem;
\item any status for $\varphi$ as a derived constant of a
  deeper theory.
\end{itemize}
The connection of $\varphi$ to downstream cascade mathematics
(in particular to the ring $\Z[\varphi]$, the icosian ring,
the phason complement, and the $\alpha$-chain) is the
subject of other sections of this paper and of the
companion $\alpha$-chain manuscripts. Those connections use
\S F1 \emph{only} at the axiomatic / named-constant level.
\end{remark}

%=====================================================================
\section{Introduction and Audit}
\label{sec:intro}
%=====================================================================

\subsection{The systemic gap}
\label{subsec:gap}

The six open Clay Millennium Prize Problems addressed in companion
manuscripts \cite{MillenniumYM, MillenniumRH, MillenniumNS,
MillenniumHodge, MillenniumBSD, MillenniumPNP} share a common
structural feature: each invokes a \emph{cascade-to-classical
projection} as a black box. Specifically:

\begin{itemize}
\item the Yang--Mills manuscript \cite{MillenniumYM} invokes a map
  $\Pi_{\gge}$ from cascade octonionic edge data to $\mathrm{SU}(3)$
  gauge connections;
\item the Riemann Hypothesis manuscript \cite{MillenniumRH} invokes
  a Hilbert--P\'olya-type operator $T_\zeta$ constructed from
  $H_4$-sector cascade spectral data;
\item the Navier--Stokes manuscript \cite{MillenniumNS} invokes a
  coarse-graining map $\Pi_{\hyd}$ from $\sigma$-fixed $D_4$-sector
  dynamics to macroscopic velocity--pressure fields;
\item the Hodge manuscript \cite{MillenniumHodge} invokes a
  ``$\sigma$-fixed part equals rational part'' projector on
  $\Q(\pphi)$-valued cohomology;
\item the Birch--Swinnerton-Dyer manuscript \cite{MillenniumBSD}
  invokes a Hecke-type transfer operator $T_E$ constructed from
  modular-symbol cascade data;
\item the $P$ vs $NP$ manuscript \cite{MillenniumPNP} invokes a
  transport between cascade Turing machines (CTMs) and ordinary
  Turing machines (TMs).
\end{itemize}

In every case, independent review at Clay-submission rigour
(\cite[iteration 1, issue \#2]{CodexReview}) has flagged the
projection itself as insufficiently defined: domains, codomains,
measures, function-space structure, and continuity properties are
not specified with enough precision for a classical referee to
verify the downstream identifications. This paper closes that
single systemic gap.

\subsection{What this paper contributes}
\label{subsec:contribution}

We construct seven projections as precise mathematical maps,
equip each with continuity/measurability/equivariance theorems,
and isolate the remaining unresolved analytic or combinatorial
issues into four explicitly named open hypothesis families. The
correspondence theorems themselves (maps well-defined, equivariant,
continuous on their natural domains) are proved unconditionally.
The coarse-graining and continuum-limit statements that would take
these projections all the way to classical continuum objects
(smooth metrics, continuum gauge fields, completed trace identities,
rational Hodge classes, asymptotic complexity separations) are
stated as conditional theorems, with the conditional hypotheses
audited honestly against the cascade-derivation sources.

More precisely, the unconditional contributions of this paper are:
\begin{enumerate}[label=(U\arabic*)]
\item Construction of the finite-level cascade configuration spaces
  $X_n^0, X_n^1, X_n^{\co}$ as separable Hilbert / vector spaces with
  explicit measures and symmetry actions (Section~\ref{sec:substrate}).
\item Construction of the inverse/cylindrical limit
  $X_\cas$ on admissible sectors, with compatibility under
  refinement, $\sigma$-involution, and rung restriction
  (Section~\ref{sec:dynamics}).
\item Finite-level cascade gradient flow: existence, uniqueness,
  energy decrease (Section~\ref{sec:dynamics}, Theorem~\ref{thm:flow-exists}).
\item Definition of the discrete moment tensor $M_n$ and of the
  finite-level metric projection $\Pi_{\met}^{(n)}$; proof that
  $\Pi_{\met}^{(n)}$ is well-defined, continuous, and
  $W(D_4)$-equivariant (Section~\ref{sec:metric-hydro},
  Theorem~\ref{thm:pimet-n}).
\item Definition of the finite-level gauge projection
  $\Pi_{\gge}^{(n)}$ from octonionic edge data to cylindrical
  $\mathrm{SU}(3)$ holonomy data, and proof that it is
  gauge-equivariant and compatible with the $\mathrm{SU}(3) \subset
  G_2 = \Der(\OO)$ identification (Section~\ref{sec:gauge},
  Theorem~\ref{thm:pigauge-n}).
\item Construction of the regularized transfer operators
  $T_\zeta[h]$ on $L^2(\R_+, dx/x)$ and $T_E[h]$ on the
  modular-symbol space, and proof that each is a densely-defined
  symmetric / bounded operator with explicit
  symmetry-intertwining properties (Section~\ref{sec:spectral},
  Theorems~\ref{thm:tzeta-h}--\ref{thm:tE-h}).
\item Proof of the $\sigma$-fixed-part-equals-rational-part theorem
  $\Fix(\sigma; H^*(X, \Q(\pphi))) = H^*(X, \Q)$, and
  functoriality of $\Pi_\sigma$ under pullback along cascade-compatible
  embeddings (Section~\ref{sec:cohomology-comp},
  Theorems~\ref{thm:fix-sigma-rational}--\ref{thm:pisigma-functorial}).
\item Construction of the CTM/TM transport maps
  $(\Pi_{\TM}, I_{\TM})$ and proof that they preserve decidability
  and polynomial time up to polynomial overhead
  (Section~\ref{sec:cohomology-comp}, Theorem~\ref{thm:ctm-tm-transport}).
\item Compatibility lemmas: refinement, $\sigma$-involution, and
  cascade dynamics all commute with the unconditional parts of the
  projections on their natural domains (Section~\ref{sec:compat}).
\item Worked examples recomputing the 24-cell moment tensor, a
  basic octonionic holonomy, a $\sigma$-fixed cohomology pair, and
  a CTM encoding of a standard TM
  (Section~\ref{sec:examples}).
\end{enumerate}

The conditional contributions of this paper are four
correspondence theorems, each under its own hypothesis stack
(Section~\ref{sec:handoff}). The stacks involve four named
hypothesis families $\mathbf{H_1}$--$\mathbf{H_4}$ \emph{and}
hypothesis imports from the hardened substrate papers P5--P7,
plus in the metric/hydrodynamic case a foundations-level
conjecture $\mathbf{H_1'}$ controlling the extension of P4's
mass-only flow-intertwining to the full closure flow:
\begin{enumerate}[label=(C\arabic*)]
\item \emph{Under $\mathbf{H_1}$} (geometric-limit hypothesis)
  \emph{together with the P5--P7 structural imports
  (Hypotheses~\ref{hyp:coord-inj-import}, \ref{hyp:tiling-import},
  \ref{hyp:separation-import}, \ref{hyp:gradient-image-import})
  and (for the full-flow case) the foundations-level conjecture
  $\mathbf{H_1'}$ (Hypothesis~\ref{hyp:full-flow})}:
  $\Pi_{\met}$ admits a continuum limit landing in $\Met(M)$ for
  a smooth compact manifold $M$, and the cascade linearized
  dynamics projects to the metric-side operator; $\Pi_{\hyd}$
  intertwines the cascade flow with the macroscopic Navier--Stokes
  operator. Well-definedness and the mass-only sub-case of the
  flow projection do \emph{not} require $\mathbf{H_1'}$; only
  the full-flow (Ricci-Laplacian-mass / NS) PDE form does.
  (Theorems~\ref{thm:pimet-cond} and~\ref{thm:pihyd-cond}.)
\item \emph{Under $\mathbf{H_2}$} (gauge-reconstruction hypothesis)
  \emph{together with the modeling choices (M1)--(M3) of
  \S\ref{subsec:gauge-modeling}} (fixed imaginary unit $i_0$,
  associator-respecting derivation lift $\iota_{\mathfrak{g}_2}$,
  and induced gauge action on octonionic edge data):
  $\Pi_{\gge}$ admits a continuum completion to $L^2_1$ classical
  connections, and transports the restricted cascade closure action
  to the continuum Yang--Mills action. (Theorem~\ref{thm:pigauge-cond}.)
\item \emph{Under $\mathbf{H_3}$} (spectral--arithmetic trace
  hypothesis): the cutoff removal $h \to \mathbf{1}$ is compatible
  with the explicit formulas for $\zeta$ and $L(E, s)$; under this
  hypothesis, the renormalized traces of $T_\zeta[h]$ and $T_E[h]$
  match the classical explicit formulas.
  (Theorem~\ref{thm:trace-match-cond}.)
\item \emph{Under $\mathbf{H_4}$} (universality / normal-form
  hypothesis): cascade embeddability is universal for smooth
  projective varieties ($\mathbf{H_{4a}}$) / cascade polynomial-time
  algorithms admit a normal form ($\mathbf{H_{4b}}$); the Hodge and
  $P$ vs $NP$ transport steps used in the companion papers become
  classical statements. (Theorem~\ref{thm:universality-cond}.)
\end{enumerate}

\subsection{Explicit audit of upstream cascade sources}
\label{subsec:audit-upstream}
\label{sec:audit}

A recurring failure mode in the earlier Millennium manuscripts is
citation of cascade-derivation results that do not in fact exist
at the cited location (``phantom citations''), or citation as
theorem of material the cited source itself labels as ``formal
framework'' or ``working hypothesis''. We classify every cascade
source used in this paper into one of three grades.

\begin{notation}[Source grades]
\label{not:grades}
We tag each cited cascade-derivation source with exactly one of
the following grades:
\begin{description}
\item[\textbf{Theorem-grade (T).}] The cited source proves a
  statement with complete mathematical content. We cite such
  results without reproof (though we may restate for clarity).
\item[\textbf{Formal framework (F).}] The cited source contains a
  statement with a proof sketch and/or numerical evidence, but
  the proof is not complete at Clay-submission rigour. We
  \emph{never} use such results as unconditional inputs to a
  theorem in the present paper. Instead, we either (a) reprove
  the specific quantitative fact we need directly from explicit
  formulas, or (b) isolate the use under one of
  $\mathbf{H_1}$--$\mathbf{H_4}$.
\item[\textbf{Motivational (M).}] The cited source provides a
  picture or heuristic but not a mathematical claim we rely on.
  We use such citations only to orient the reader; no theorem
  in this paper depends on any motivational citation.
\end{description}
\end{notation}

The audit of cascade sources used in the present paper is given in
Table~\ref{tab:audit}. Sources not listed in the table are not used.

\begin{table}[h]
\caption{Audit of upstream cascade-derivation sources used in this
paper.}
\label{tab:audit}
\begin{tabular}{llll}
\hline
Source & Section(s) used & Grade & Role in present paper \\
\hline
\cite{CascadeFoundations} \S F1 & Sec.~\ref{sec:substrate} & T &
  $\pphi$-uniqueness from self-similarity \\
\cite{CascadeFoundations} \S F2 & Sec.~\ref{sec:dynamics} & T &
  closure-functional form $\alpha R + \beta E - \gamma Q$ \\
\cite{CascadeFoundations} \S F3 & Sec.~\ref{sec:substrate} & T &
  seven-rung cascade from Coxeter classification \\
\cite{CascadeFoundations} \S F5 & Sec.~\ref{sec:dynamics} & T &
  $\sigma$-involution on $\Q(\pphi)$-coefficients \\
\cite{CascadeEmbeddings} \S\S 2--3 & Sec.~\ref{sec:substrate} & T &
  $D_4 \subset E_8$ inclusion; $H_4$ icosian projection \\
\cite{CascadeEmbeddings} \S 5 & Sec.~\ref{sec:substrate} & T &
  $4$-cube and $\mathrm{Cl}(1,3)$ coincidence at rung 16 \\
\cite{CascadeEmbeddings} \S 6 & Sec.~\ref{sec:substrate},
  \ref{sec:gauge} & T & $E_8$ self-dimension as rung 8; octonionic
  structure \\
\cite{CascadeGR} \S C4.1 & Sec.~\ref{sec:metric-hydro} & T &
  discrete moment tensor formula for point source on 24-cell \\
\cite{CascadeGR} \S C2.bis & Sec.~\ref{sec:metric-hydro},
  \ref{sec:handoff} & F & GH continuum limit (kept under $\mathbf{H_1}$) \\
\cite{CascadeGR} \S C3.bis & Sec.~\ref{sec:metric-hydro} & F &
  $A_5$-density in $\mathrm{SO}(3)$ (used only as motivation for
  $\mathbf{H_1}$) \\
\cite{CascadeDynamics} \S\S D1--D4 & Sec.~\ref{sec:dynamics},
  \ref{sec:cohomology-comp} & T & primitive vocabulary; CTM
  definitions \\
\cite{CascadeDynamics} \S\S D5--D10 & Sec.~\ref{sec:cohomology-comp} & F &
  cascade complexity classes (kept under $\mathbf{H_{4b}}$) \\
\cite{CascadeMillennium} \S\S MP1--MP7 & Sec.~\ref{sec:handoff} & M &
  orientation to the Millennium picture \\
\hline
\end{tabular}
\end{table}

Sources excluded from this paper's load-bearing machinery are
listed in Table~\ref{tab:exclusions}; they appear in earlier draft
Millennium manuscripts but are not used here. This exclusion is
the principal reason the present paper is narrower, and
correspondingly more reliable, than those drafts.

\begin{table}[h]
\caption{Cascade-derivation sources excluded from load-bearing use
in this paper.}
\label{tab:exclusions}
\begin{tabular}{llp{2.8in}}
\hline
Source & Reason & Handling \\
\hline
\cite{CascadePregeometry} \S\S P4--P8 & Inconsistent internal
  notation; not self-contained at theorem level; referenced
  ``VFD Master Math'' manuscript not present in workspace &
  Use only as motivational (M) background for the $\pphi$-entropy
  functional; do not cite as theorem input \\
\cite{CascadeFoundations} \S\S F6--F8 & F6 claims Burago--Ivanov
  hypotheses verified --- but the verification in the cited
  source does not uniformly bound the $g/\mathrm{geo}$ overshoot
  flagged in \cite[\S C2.bis]{CascadeGR}; F8 is
  closure-coefficient matching, which is circular as an upstream
  input for a \emph{correspondence} paper whose subject is
  building the projections on which closure is computed &
  F6: replaced by explicit open item in $\mathbf{H_1}$. F8: not
  cited. F7 used only for orientation, not as theorem input \\
\cite{CascadeMillennium} \S\S MP2, MP4, MP5, MP7 cascade-internal
  ``proof sketches'' & Each contains technical items the source
  itself flags as remaining work & Cited in
  Section~\ref{sec:handoff} only as target for the handoff table;
  no theorem in this paper depends on them \\
\hline
\end{tabular}
\end{table}

\begin{remark}[On the choice to exclude rather than upgrade]
\label{rmk:exclude-not-upgrade}
The tables above reflect a deliberate choice: where a
cascade-derivation source is labelled in its own narrative as
``formal framework'' or contains open technical items flagged by
the source author, we do not attempt to upgrade that source within
the present paper. Upgrading is the separate task performed either
(i) by completing the open items in the source itself, or (ii)
classically, e.g.\ via Burago--Ivanov
\cite{BuragoBuragoIvanov} / Cheeger--Colding
\cite{CheegerColding1997}. Until either route is executed, we keep
the corresponding statement as part of the appropriate
$\mathbf{H}_i$ hypothesis.
\end{remark}

\subsection{Relation to the six Millennium papers}
\label{subsec:relation-to-millennium}

The six Millennium manuscripts will, after the present paper is in
place, cite specific definitions and theorems here to close their
projection-definition gap. The handoff table is given in
Section~\ref{sec:handoff} (Table~\ref{tab:handoff}); what each
Millennium paper inherits and what remains open after that
inheritance is stated there precisely. A summary: each of the
five \emph{constructive} target classes
(YM, RH, NS, BSD, $P$ vs $NP$) inherits exactly one of
$\mathbf{H_1}$--$\mathbf{H_4}$, and the Hodge paper inherits
$\mathbf{H_{4a}}$.

The Poincar\'e manuscript \cite{MillenniumPoincare} is a special
case: Perelman's theorem closes Poincar\'e at Clay bar already,
and the cascade contribution per
\cite[\S MP6.3]{CascadeMillennium} is explicitly ``no contribution
beyond consistency''. The present paper reconciles the
consistency claim made there with the precise metric-projection
machinery $\Pi_{\met}$ constructed here.

\subsection{Organization and notational conventions}
\label{subsec:organization}

Section~\ref{sec:substrate} constructs the discrete substrate:
refinement graphs, counting measures, sector fibres, and
finite-level sector spaces, culminating in Theorem~\ref{thm:hilbert}
(finite-level spaces are separable Hilbert; refinement maps are
measurable). Section~\ref{sec:dynamics} defines configuration
spaces at inverse/cylindrical limit, closure-energy functionals,
the finite-level gradient flow, and the symmetry actions;
Theorems~\ref{thm:flow-exists} and \ref{lem:commute} establish
existence / uniqueness / commutation. Section~\ref{sec:corresp}
sets the abstract correspondence scheme and proves a general
measurability / equivariance lemma. Sections~\ref{sec:metric-hydro},
\ref{sec:gauge}, \ref{sec:spectral}, and \ref{sec:cohomology-comp}
construct the four families of projections:
metric/hydrodynamic, gauge, spectral-arithmetic, and
cohomological / computational. Section~\ref{sec:compat} establishes
the compatibility diagrams. Section~\ref{sec:examples} gives the
worked examples. Section~\ref{sec:handoff} states the conditional
correspondence theorems and the handoff table to the six Millennium
papers. An appendix audits notation and citations.

\begin{notation}[Global conventions]
\label{not:conventions}
Throughout the paper:
\begin{itemize}
\item $\pphi = (1 + \sqrt 5)/2$ denotes the golden ratio; $\bar\pphi
  = 1 - \pphi = -1/\pphi$ is its Galois conjugate. The notation
  $\Q(\pphi) := \Q[\pphi]$ denotes the quadratic field $\Q[\sqrt 5]$.
\item $\sigma \colon \Q(\pphi) \to \Q(\pphi)$ is the non-trivial
  Galois involution fixing $\Q$: $\sigma(\pphi) = \bar\pphi$,
  $\sigma(a + b\pphi) = a + b\bar\pphi$. We extend $\sigma$
  coefficientwise to all $\Q(\pphi)$-modules.
\item $H_4$ denotes the non-crystallographic Coxeter group of type
  $H_4$ acting on $\R^4$ as the full symmetry group of the 600-cell
  / 120-cell dual pair \cite[\S 3]{CascadeEmbeddings},
  \cite{CoxeterRegularPolytopes}. $|H_4| = 14\,400$.
\item $W(D_4) \subset W(B_4) \subset \mathrm{O}(4)$ is the Weyl
  group of type $D_4$, order $192$. Triality is the outer
  automorphism of $D_4$ of order $3$, preserving the $24$-cell
  \cite[\S 2]{CascadeEmbeddings}.
\item $\OO$ denotes the normed division algebra of octonions;
  $\Im(\OO) \cong \R^7$ is the imaginary part; $G_2 = \Der(\OO)$
  is its derivation algebra; $\mathrm{SU}(3) \subset G_2$ is the
  stabilizer of any fixed imaginary unit \cite[\S 6]{CascadeEmbeddings}.
\item Graphs are undirected, simple, and locally finite unless
  otherwise stated. $V(G)$ denotes vertices, $E(G)$ unoriented
  edges, $E^{\rm or}(G)$ oriented edges.
\item All function spaces $L^p$ use Lebesgue measure on $\R^d$
  unless otherwise specified; $\ell^2(V, \mu; V_{\rm fib})$ denotes
  square-summable functions from the measurable set $V$ with
  measure $\mu$ into the fibre Hilbert space $V_{\rm fib}$, with
  inner product $\langle f, g \rangle = \sum_{v \in V} \mu(\{v\})
  \langle f(v), g(v) \rangle_{V_{\rm fib}}$.
\item $\Pi_*^{(n)}$ denotes the finite-level (refinement level $n$)
  projection, and $\Pi_*$ denotes the corresponding inverse-limit /
  continuum limit (which may be conditional on one of the
  $\mathbf{H}_i$).
\end{itemize}
\end{notation}

The reader should keep the audit of Table~\ref{tab:audit} in mind
throughout. In particular, the word ``cascade'' in this paper
refers only to the theorem-grade cascade inputs listed there; the
wider cascade literature contains additional material which we
deliberately do not use.

%=====================================================================
\section{The Discrete Cascade Substrate}
\label{sec:substrate}
%=====================================================================

This section constructs the discrete substrate on which all
cascade data lives at a given refinement level. The construction
is entirely explicit: graphs, counting measures, fibre spaces,
and $\ell^2$-sections. The main result (Theorem~\ref{thm:hilbert})
is that the resulting sector spaces form a compatible family of
separable Hilbert (or $\Q(\pphi)$-) vector spaces.

\subsection{Refinement graphs}
\label{subsec:graphs}

\begin{definition}[The base graphs $G_0^{H_4}$ and $G_0^{D_4}$]
\label{def:base-graphs}
Let $G_0^{H_4}$ denote the $1$-skeleton of the regular 600-cell in
$\R^4$: its vertex set is the $120$-element vertex set of the
600-cell (the icosians of norm $1$, see
\cite[\S 3]{CascadeEmbeddings}), and two vertices are joined by an
undirected edge iff their Euclidean distance equals the edge
length of the 600-cell. Let $G_0^{D_4}$ denote the $1$-skeleton of
the regular $24$-cell inscribed in the 600-cell via the standard
$D_4 \subset H_4$ sub-root-system inclusion
\cite[\S 2]{CascadeEmbeddings}: its vertex set is the $24$-element
set of $D_4$ minuscule vectors, and edges join pairs of vertices at
the $24$-cell edge distance.
\end{definition}

\begin{definition}[Schl\"afli refinement]
\label{def:schlafli}
The \emph{Schl\"afli refinement} of a polytope graph $G$, denoted
$\mathrm{Sch}(G)$, replaces each edge $e = \{u, v\}$ by the unique
regular edge subdivision respecting the Coxeter symmetry of $G$, and
adjoins new vertices at the face centroids in accordance with the
$2I/2T$ coset decomposition of the binary icosahedral / binary
tetrahedral groups as described in
\cite[\S C2.bis.1]{CascadeGR}. We do not require the full
convergence claim of \cite[\S C2.bis]{CascadeGR} (which is formal
framework, grade F); we use only the combinatorial definition of
the refinement map itself, which is algorithmic.

Let $G_{n+1}^{\bullet} := \mathrm{Sch}(G_n^{\bullet})$ for
$\bullet \in \{H_4, D_4\}$. At each level, the resulting graph is
simple, locally finite, and carries a free action of the relevant
Coxeter symmetry group ($W(H_4)$ or $W(D_4)$).
\end{definition}

\begin{remark}[Status of the refinement used here]
\label{rmk:refinement-status}
Only the existence of the refinement operation and its
combinatorial structure are used below (these are
theorem-grade per Definition~\ref{def:schlafli}). The geometric
claim that the sequence $(G_n^{H_4}, d_n) \to (S^3, d_{\rm round})$
in Gromov--Hausdorff sense is \emph{not} used as an unconditional
input. Whenever we wish to invoke such a continuum limit, we do so
under $\mathbf{H_1}$; see Sections~\ref{sec:metric-hydro}
and \ref{sec:handoff}.
\end{remark}

\subsection{Counting measures}
\label{subsec:counting-measures}

\begin{definition}[Counting measures $\mu_n$ and $\nu_n$]
\label{def:measures}
Let $G_n$ be any of $G_n^{H_4}, G_n^{D_4}$. Define:
\begin{align*}
  \mu_n(S) &:= \frac{|S|}{|V(G_n)|}, \quad S \subseteq V(G_n);
  \\
  \nu_n(T) &:= \frac{|T|}{|E^{\rm or}(G_n)|}, \quad T \subseteq E^{\rm or}(G_n);
\end{align*}
where $E^{\rm or}(G_n)$ is the set of \emph{oriented} edges
(each unoriented edge $\{u, v\}$ contributing two oriented edges
$(u, v)$ and $(v, u)$). Both $\mu_n$ and $\nu_n$ are probability
measures on their respective finite sets.
\end{definition}

\begin{remark}[Normalization]
Normalization to probability measure is a convenience: it ensures
that the $\ell^2$ inner product
$\langle f, g\rangle_{\mu_n} = \sum_v \mu_n(\{v\}) \langle f(v),
g(v)\rangle_{V_{\rm fib}}$ is bounded independently of $n$ when
$f, g$ are uniformly bounded. Alternative normalizations (e.g.\
counting measure scaled by $\varepsilon$-net radius $h_n^d$) lead
to equivalent sector spaces up to a refinement-level-dependent
rescaling; we use probability normalization throughout.
\end{remark}

\subsection{Sector fibres}
\label{subsec:fibres}

The cascade seven-rung structure
$E_8 \to H_4 \to 40 \to D_4 \to 16 \to 8 \to 0$ is established in
\cite[\S F3]{CascadeFoundations}. Each rung carries a canonical
vector-space fibre, which we recall from
\cite[\S\S 2--6]{CascadeEmbeddings}:

\begin{definition}[Sector fibres]
\label{def:fibres}
Define the sector fibres as follows (with stated dimensions over $\R$
and, where relevant, $\Q(\pphi)$-module structure):
\begin{enumerate}
\item $V_{H_4} := \R^4$ with the standard $H_4$-module structure
  given by the reflection representation. Over $\Q(\pphi)$, the
  integer-lattice icosian model
  \cite[\S 3.1]{CascadeEmbeddings} gives $V_{H_4}$ a natural
  $\Q(\pphi)^4$-module structure; we use the $\R^4$ model by
  default and note the $\Q(\pphi)^4$ model when $\sigma$-actions
  are needed.
\item $V_{D_4} := \R^4$ with the standard $D_4$-reflection
  representation; this is the representation on the $24$-cell
  vertex span \cite[\S 2]{CascadeEmbeddings}.
\item $V_{16} := \mathrm{Cl}(1,3)$, the real Clifford algebra of
  signature $(1, 3)$, of real dimension $16$
  \cite[\S 5.1]{CascadeEmbeddings}.
\item $V_{8} := \OO$, the octonion algebra, of real dimension
  $8$ \cite[\S 6]{CascadeEmbeddings}. The imaginary subspace
  $\Im(\OO)$ has real dimension $7$ and will host gauge data
  (Section~\ref{sec:gauge}).
\end{enumerate}
Each $V_\bullet$ carries the inner product induced by restriction
from $E_8 \otimes \R \cong \R^8$ via the respective embedding
\cite[\S\S 2--6]{CascadeEmbeddings}; with this inner product each
$V_\bullet$ is a finite-dimensional Euclidean space.
\end{definition}

\begin{remark}[Why these four and not others]
The choice of fibres $V_{H_4}, V_{D_4}, V_{16}, V_8$ matches the
rungs carrying non-trivial cascade dynamics for the Millennium
applications: $V_{H_4}$ for spectral data (RH, BSD through the
modular-symbol map), $V_{D_4}$ for metric and hydrodynamic data
(Poincar\'e, NS), $V_{16}$ for the $\mathrm{Cl}(1,3)$ structure
underlying the real geometric-algebra formulation, and $V_8$ for
gauge data (YM). The remaining rungs (40 and 0) do not carry
additional fibre content beyond what is already present in
$V_{H_4}$ and $V_{D_4}$; we omit them.
\end{remark}

\subsection{Finite-level sector spaces}
\label{subsec:sector-spaces}

\begin{definition}[Vertex-sector space $X_n^0$]
\label{def:X-n-0}
For $n \geq 0$ let
\[
  X_n^0 := \ell^2\big(V(G_n), \mu_n; V_{H_4} \oplus V_{D_4}
    \oplus V_{16} \oplus V_8\big)
\]
with inner product
\[
  \langle f, g\rangle_{X_n^0}
  := \sum_{v \in V(G_n)} \mu_n(\{v\})
    \big\langle f(v), g(v) \big\rangle_{V_{H_4} \oplus V_{D_4}
    \oplus V_{16} \oplus V_8}.
\]
Here $V(G_n)$ may denote either $V(G_n^{H_4})$ or $V(G_n^{D_4})$;
we suppress the decoration when the choice is fixed by context.
\end{definition}

\begin{definition}[Edge-sector space $X_n^1$]
\label{def:X-n-1}
For $n \geq 0$ let
\[
  X_n^1 := \ell^2\big(E^{\rm or}(G_n), \nu_n; \Im(\OO)\big)
\]
with inner product
\[
  \langle A, B\rangle_{X_n^1}
  := \sum_{e \in E^{\rm or}(G_n)} \nu_n(\{e\})
    \big\langle A(e), B(e) \big\rangle_{\Im(\OO)},
\]
subject to the oriented-edge reversal convention $A(-e) = -A(e)$
for every $e \in E^{\rm or}(G_n)$, where $-e$ denotes the edge $e$
with reversed orientation. This is the standard gauge-theoretic
convention for $1$-cochains.
\end{definition}

\begin{definition}[Cohomology-sector space $X^{\co}(X)$]
\label{def:X-coh}
For any smooth projective complex variety $X$ (more generally, any
finite-dimensional $\Q(\pphi)$-vector-space-valued cohomology
theory satisfying the axioms of Section~\ref{sec:cohomology-comp}),
let
\[
  X^{\co}(X) := H^*\big(X, \Q(\pphi)\big)
  = H^*(X, \Q) \otimes_{\Q} \Q(\pphi).
\]
This is a finite-dimensional $\Q(\pphi)$-vector space with a
natural $\sigma$-action (coefficientwise), and its $\sigma$-fixed
part is $H^*(X, \Q)$; this fact is established as
Theorem~\ref{thm:fix-sigma-rational}.
\end{definition}

\subsection{The Hilbert / vector space structure theorem}
\label{subsec:hilbert-thm}

\begin{theorem}[Sector spaces: basic structural properties]
\label{thm:hilbert}
Fix $n \geq 0$ and $\bullet \in \{H_4, D_4\}$.
\begin{enumerate}
\item $X_n^0$ and $X_n^1$ are finite-dimensional real Hilbert
  spaces, hence in particular separable Hilbert spaces. Their
  dimensions are
  \[
    \dim_\R X_n^0 = |V(G_n^\bullet)| \cdot
    (\dim V_{H_4} + \dim V_{D_4} + \dim V_{16} + \dim V_8)
    = |V(G_n^\bullet)| \cdot 32,
  \]
  and
  \[
    \dim_\R X_n^1 = \tfrac{1}{2} |E^{\rm or}(G_n^\bullet)| \cdot
    \dim \Im(\OO) = \tfrac{1}{2} |E^{\rm or}(G_n^\bullet)| \cdot 7,
  \]
  where the factor $\tfrac{1}{2}$ in the edge-sector dimension
  reflects the reversal convention (each unoriented edge
  contributes $7$ real degrees of freedom).
\item For any smooth projective complex variety $X$, $X^{\co}(X)$
  is a finite-dimensional $\Q(\pphi)$-vector space. Its dimension
  over $\Q$ is $2 \dim_{\Q} H^*(X, \Q)$.
\item For each $n \geq 0$, the refinement map
  $\mathrm{Sch}_\ast \colon G_n^\bullet \to G_{n+1}^\bullet$
  (inclusion of the level-$n$ vertex set into the level-$(n+1)$
  vertex set, and the corresponding map on oriented edges) induces
  bounded linear maps
  $R_n^0 \colon X_n^0 \to X_{n+1}^0$ and
  $R_n^1 \colon X_n^1 \to X_{n+1}^1$
  defined by zero-extension outside the image of
  $\mathrm{Sch}_\ast$; these maps are Borel measurable and
  $W(G_n^\bullet)$-equivariant for the Coxeter-symmetry group
  of $G_n^\bullet$.
\end{enumerate}
\end{theorem}

\begin{proof}
\emph{(1)} For each $v \in V(G_n^\bullet)$, the fibre
$V_{H_4} \oplus V_{D_4} \oplus V_{16} \oplus V_8$ is a
finite-dimensional Euclidean space (by Definition~\ref{def:fibres})
of real dimension $4 + 4 + 16 + 8 = 32$. The vertex set
$V(G_n^\bullet)$ is finite (explicitly: the Schl\"afli refinement
multiplies the vertex count by a fixed factor depending on the
polytope type, but remains finite at each level), so $X_n^0$ is a
finite direct sum of copies of this fibre space, weighted by the
probability measure $\mu_n$. A finite direct sum of Euclidean
spaces with positive weights is Hilbert. The dimension count is
immediate.

The edge-sector analysis is identical: $\Im(\OO) \cong \R^7$
(Definition~\ref{def:fibres}), $E^{\rm or}(G_n^\bullet)$ is finite,
and the reversal convention
$A(-e) = -A(e)$ cuts the effective real dimension by a factor
of two.

\emph{(2)} $H^*(X, \Q)$ is finite-dimensional by classical Hodge
theory \cite{GriffithsHarris}. Tensoring with $\Q(\pphi) = \Q[\pphi]$,
a $\Q$-vector space of dimension $2$, doubles the $\Q$-dimension
but is naturally a $\Q(\pphi)$-module of the same rank as
$H^*(X, \Q)$ has $\Q$-dimension.

\emph{(3)} The refinement map on vertices is, by
Definition~\ref{def:schlafli}, a combinatorial embedding
$V(G_n^\bullet) \hookrightarrow V(G_{n+1}^\bullet)$. Extending a
function $f \in X_n^0$ by zero outside this embedding defines
$R_n^0 f \in X_{n+1}^0$ pointwise. Boundedness: for any $f \in X_n^0$,
\[
  \|R_n^0 f\|_{X_{n+1}^0}^2
  = \sum_{v \in V(G_{n+1}^\bullet)} \mu_{n+1}(\{v\})
    \|(R_n^0 f)(v)\|^2
  = \sum_{v \in V(G_n^\bullet)} \mu_{n+1}(\{v\}) \|f(v)\|^2
  \leq C_n \|f\|_{X_n^0}^2,
\]
where $C_n := \max_{v \in V(G_n^\bullet)}
\mu_{n+1}(\{v\})/\mu_n(\{v\}) =
|V(G_n^\bullet)|/|V(G_{n+1}^\bullet)|$ is a finite constant
independent of $f$. Borel measurability on finite-dimensional spaces
is automatic for linear maps. $W$-equivariance follows from the
Coxeter-symmetry compatibility of the Schl\"afli refinement
recorded in \cite[\S 12.2]{CascadeEmbeddings} (this is a
theorem-grade input: explicit vertex coordinates are given there).
The edge-sector proof is identical.
\end{proof}

\begin{remark}[Explicit dimensions at low levels]
\label{rmk:low-dims}
For orientation: $|V(G_0^{H_4})| = 120$ (600-cell vertices), so
$\dim_\R X_0^{0,H_4} = 120 \cdot 32 = 3{,}840$; $|V(G_0^{D_4})| =
24$, so $\dim_\R X_0^{0,D_4} = 768$.
$|E^{\rm or}(G_0^{H_4})| = 720 \cdot 2 = 1{,}440$ (the 600-cell has
720 edges, each oriented both ways), so
$\dim_\R X_0^{1,H_4} = 720 \cdot 7 = 5{,}040$, etc.
At level $n$, all dimensions grow polynomially or exponentially in
$n$ depending on the refinement rule, but remain finite at each
level.
\end{remark}

\subsection{Inverse-limit structure and compatibility}
\label{subsec:inverse-limit}

The finite-level sector spaces form a directed system under the
refinement maps $R_n^\bullet$. We now record the inverse-limit
(equivalently cylindrical) structure, which will support the
infinite-level constructions of Sections~\ref{sec:dynamics}
and beyond.

\begin{definition}[Cylindrical test-vectors and the inverse limit]
\label{def:inverse-limit}
A \emph{cylindrical} vector $\phi$ of degree $n$ in the vertex
sector is, by definition, a vector $\phi_n \in X_n^0$. Two
cylindrical vectors $(n, \phi_n)$ and $(n', \phi_{n'})$ with
$n \leq n'$ are declared \emph{equivalent} if
$R_{n'-1}^0 \circ \cdots \circ R_n^0 \phi_n = \phi_{n'}$. The
\emph{cylindrical sector space}
\[
  X^0_\infty := \bigcup_{n \geq 0} X_n^0 / \sim
\]
is the set of equivalence classes; it carries a natural pre-Hilbert
structure by using the level-$n$ inner product on any
representative. The \emph{inverse-limit sector space}
$X^0_{\varprojlim}$ is the completion of $X^0_\infty$ with respect
to the pre-Hilbert norm; it is a separable Hilbert space by
standard inverse-limit / projective-system theory
\cite[\S 6.2]{ReedSimonI}.

An analogous construction gives the edge-sector inverse limit
$X^1_{\varprojlim}$.
\end{definition}

\begin{remark}[Role of the inverse limit]
The inverse-limit construction is the standard device for
passing from finite-level discrete data to an idealized continuum
object \cite{ReedSimonI}. We do \emph{not} claim here that
$X^0_{\varprojlim}$ is isomorphic to any specific classical
function space (e.g.\ $L^2(S^3)$ with its round-metric measure); that
identification depends on $\mathbf{H_1}$ and is taken up in
Section~\ref{sec:metric-hydro}. The inverse-limit structure is
\emph{combinatorial} here; the geometric interpretation is
downstream.
\end{remark}

\subsection{Symmetry actions on the substrate}
\label{subsec:symmetries}

We record the symmetry actions that the finite-level sector spaces
support. All actions commute with the $\ell^2$ structure and with
zero-extension, so they descend to $X^0_\infty$ and to the
inverse-limit $X^0_{\varprojlim}$.

\begin{definition}[Coxeter actions]
\label{def:coxeter-actions}
The Coxeter group $W(H_4)$ acts on $V(G_n^{H_4})$ by its standard
action on the 600-cell / its refinements
\cite[\S 3]{CascadeEmbeddings}, and correspondingly on
$V_{H_4} \subset \R^4$ through its reflection representation. These
two actions define a simultaneous action of $W(H_4)$ on
$X_n^{0,H_4}$ by
\[
  (w \cdot f)(v) := \rho(w) \, f(w^{-1} \cdot v)
  \quad \text{for } w \in W(H_4), v \in V(G_n^{H_4}),
\]
where $\rho$ is the reflection representation on
$V_{H_4} \oplus V_{D_4} \oplus V_{16} \oplus V_8$ (with $W(D_4)$
acting through its inclusion $W(D_4) \subset W(H_4)$, and
$W(H_4)$ acting trivially on $V_{16}$ and $V_8$ for now; the
non-trivial gauge actions enter in Section~\ref{sec:gauge}).
The action is unitary with respect to
$\langle\cdot, \cdot\rangle_{X_n^0}$.

The analogous $W(D_4)$-action on $X_n^{0, D_4}$ is defined
identically.
\end{definition}

\begin{definition}[Galois action $\sigma$]
\label{def:sigma-action-subst}
Let $\Q(\pphi)$ act on $V_{H_4}$ through the integer-lattice icosian
model \cite[\S 3.1]{CascadeEmbeddings}; let $\sigma$ act on
$V_{H_4}$ by its coefficientwise Galois action
$\sigma \colon a + b \pphi \mapsto a + b \bar\pphi$. This action
extends uniquely to an $\R$-linear involution on the $\R$-span of
$V_{H_4}$, commuting with the $H_4$-reflection representation (see
\cite[\S F5]{CascadeFoundations}, theorem-grade).

Extend $\sigma$ to $V_{D_4}, V_{16}, V_8$ by the trivial action:
these fibres carry no $\Q(\pphi)$-structure beyond the rational
structure inherited from $\Q \subset \Q(\pphi)$, so their
coefficientwise $\sigma$ is the identity.

Extend $\sigma$ to $X_n^0$ pointwise:
$(\sigma f)(v) := \sigma(f(v))$. The resulting map is $\R$-linear
and a unitary involution on $X_n^0$.
\end{definition}

\begin{remark}[On the scope of $\sigma$]
The $\sigma$-action is non-trivial on the $V_{H_4}$-sector of
$X_n^0$ and trivial on the other sectors. This is the minimal
non-trivial Galois structure consistent with
\cite[\S F5]{CascadeFoundations}. The spectral-arithmetic
consequences (Section~\ref{sec:spectral}) use only the $V_{H_4}$
action.
\end{remark}

\begin{definition}[Refinement shift]
\label{def:refine-shift}
The refinement maps $R_n^0, R_n^1$ of
Theorem~\ref{thm:hilbert}(3) constitute the \emph{refinement shift}.
Refinement shift commutes with $\sigma$ (since $\sigma$ acts
coefficientwise and zero-extension is $\Q$-linear) and, after
restriction to $X_n^{0,D_4}$, commutes with $W(D_4)$; for
$X_n^{0,H_4}$ and $W(H_4)$ the analogous statement holds.
\end{definition}

\begin{lemma}[Commutation of symmetries]
\label{lem:commute-subst}
On $X_n^0$ and $X_n^1$, the three actions (Coxeter, $\sigma$,
refinement shift) commute in all pairs.
\end{lemma}

\begin{proof}
$\sigma$ and the Coxeter actions: $\sigma$ acts coefficientwise on
$\Q(\pphi)$-coefficients; the Coxeter reflection representation is
defined over $\Q(\pphi)$ (the $H_4$ Coxeter matrix is rational with
$\pphi$-entries in the standard basis,
\cite[\S 3]{CascadeEmbeddings}), so its matrix coefficients are
$\sigma$-stable. Therefore $\sigma$ and any Coxeter element commute
at the level of matrix multiplication.

$\sigma$ and refinement: zero-extension is $\Q$-linear and
$\sigma$ is $\Q$-linear and $\R$-linear; hence they commute.

Coxeter and refinement: the Schl\"afli refinement is defined as
$W(G_n)$-equivariant by \cite[\S 12.2]{CascadeEmbeddings}
(theorem-grade); the vertex extension map is then automatically
equivariant.
\end{proof}

\begin{remark}[What Section~\ref{sec:substrate} has and has not done]
At this stage the substrate is fully defined at finite level: we
have graphs, measures, fibres, sector spaces, refinement maps,
and three commuting symmetry actions. All constructions are
entirely \emph{combinatorial and finite-dimensional}; nothing yet
depends on the geometric continuum limit (which requires
$\mathbf{H_1}$). Section~\ref{sec:dynamics} adds the closure
functional and gradient flow on top of this substrate, still at
finite level; Section~\ref{sec:corresp} then abstracts the
correspondence scheme; the sector-specific projections
$\Pi_{\met}, \Pi_{\hyd}, \Pi_{\gge}, T_\zeta[h], T_E[h], \Pi_\sigma,
\Pi_\TM$ are constructed in Sections~\ref{sec:metric-hydro} through
\ref{sec:cohomology-comp}.
\end{remark}

%=====================================================================
\section{Configuration Spaces and Dynamics}
\label{sec:dynamics}
%=====================================================================

This section equips the finite-level sector spaces with a closure
energy functional, a gradient flow, and the full inverse-limit
configuration space. All constructions remain finite-dimensional at
each refinement level; the only infinite-dimensional object is the
inverse-limit $X_\cas$, which is a separable Hilbert space by
standard projective-system theory.

\subsection{The cascade configuration space}
\label{subsec:config-space}

\begin{definition}[Cascade configuration space]
\label{def:xcas}
The \emph{cascade configuration space} is the inverse-limit Hilbert
space
\[
  X_\cas := X^0_{\varprojlim} \oplus X^1_{\varprojlim},
\]
where $X^0_{\varprojlim}$ and $X^1_{\varprojlim}$ are as in
Definition~\ref{def:inverse-limit}. We write a general element as
$\Phi = (\Phi^0, \Phi^1)$ with $\Phi^0 \in X^0_{\varprojlim}$ and
$\Phi^1 \in X^1_{\varprojlim}$, and we continue to write
$\Phi_n = (\Phi_n^0, \Phi_n^1) \in X_n^0 \oplus X_n^1$ for its
level-$n$ component (the equivalence-class representative at
level $n$).

The \emph{admissible subspace} $X_\cas^{\rm adm} \subset X_\cas$
consists of those $\Phi \in X_\cas$ whose level-$n$ components
satisfy the uniform bound
\begin{equation}
\label{eq:admissible-bound}
  \sup_{n \geq 0} \|\Phi_n\|_{X_n^0 \oplus X_n^1} < \infty.
\end{equation}
The admissible subspace is a Hilbert subspace of $X_\cas$ (it is the
inverse limit of unit balls at each level).
\end{definition}

\begin{remark}[Why admissibility matters]
The uniform bound \eqref{eq:admissible-bound} is the standard
cylindrical-measurability condition in projective-system
constructions. Without it, the finite-level gradient flow
(Section~\ref{subsec:gradient-flow}) can escape to infinity in
norm as $n \to \infty$; with it, the flow is well-defined on the
whole inverse-limit. All physical cascade configurations discussed
in the literature \cite[\S 4]{CascadeEmbeddings},
\cite[\S\S D1--D4]{CascadeDynamics} are admissible in this sense.
\end{remark}

\subsection{Closure energies}
\label{subsec:closure-energies}

The cascade closure functional of \cite[\S F2]{CascadeFoundations}
(theorem-grade, Table~\ref{tab:audit}) takes the form
\[
  \mathcal{F} = \alpha R + \beta E - \gamma Q
\]
with $\alpha, \beta, \gamma > 0$ and $R, E, Q$ the
\emph{curvature}, \emph{kinetic}, and \emph{correction} densities.
On the substrate of Section~\ref{sec:substrate}, these densities
have canonical finite-level discretizations, which we now record.
All three densities are quadratic forms on the sector space.

\begin{definition}[Finite-level closure energy $F_n$]
\label{def:Fn}
Let $\Phi = (\Phi^0, \Phi^1) \in X_n^0 \oplus X_n^1$. The
finite-level closure energy is
\begin{equation}
\label{eq:Fn}
  F_n(\Phi) := \alpha R_n(\Phi) + \beta E_n(\Phi) - \gamma Q_n(\Phi),
\end{equation}
where:
\begin{itemize}
\item $R_n$ is the \emph{discrete curvature}:
\[
  R_n(\Phi) := \tfrac{1}{2}
    \big\langle \Phi^0, \Delta_n \Phi^0 \big\rangle_{X_n^0}
    + \tfrac{1}{2}
    \big\langle \Phi^1, \partial_n^* \partial_n \Phi^1 \big\rangle_{X_n^1},
\]
with $\Delta_n$ the graph Laplacian on $X_n^0$ and
$\partial_n^* \partial_n$ the discrete exterior-derivative Laplacian
on $X_n^1$;
\item $E_n$ is the \emph{discrete kinetic}:
\[
  E_n(\Phi) := \tfrac{1}{2}
    \big\langle \mathrm{d}_n \Phi^0, \mathrm{d}_n \Phi^0 \big\rangle_{X_n^1},
\]
with $\mathrm{d}_n \colon X_n^0 \to X_n^1$ the graph-exterior
derivative (oriented-edge difference);
\item $Q_n$ is the \emph{discrete correction}:
\[
  Q_n(\Phi) := \tfrac{1}{2}
    \big\langle \Phi^0, \Phi^0 \big\rangle_{X_n^0}
    + \tfrac{1}{2}
    \big\langle \Phi^1, \Phi^1 \big\rangle_{X_n^1}.
\]
\end{itemize}
Each of $R_n, E_n, Q_n$ is a continuous, $W$-invariant, and
$\sigma$-invariant quadratic form on $X_n^0 \oplus X_n^1$.
\end{definition}

\begin{remark}[On the choice of discrete Laplacians]
The discrete Laplacians $\Delta_n$ and $\partial_n^* \partial_n$
are the standard combinatorial Laplacians associated to
$G_n$ with the measures $\mu_n, \nu_n$ of
Definition~\ref{def:measures}. These are well-studied
finite-dimensional symmetric operators; in particular they are
positive semidefinite and their spectra are finite sets with
multiplicities. Use of the combinatorial Laplacian at this
finite-level stage should not be confused with the much stronger
spectral-convergence claim of
\cite[\S C2.bis.1]{CascadeGR}, which is formal framework and which
we invoke only under $\mathbf{H_1}$.
\end{remark}

\begin{lemma}[Closure-energy regularity]
\label{lem:Fn-regularity}
The functional $F_n \colon X_n^0 \oplus X_n^1 \to \R$ is
continuously differentiable with gradient
\[
  \nabla F_n(\Phi) = \alpha \nabla R_n(\Phi) + \beta \nabla E_n(\Phi)
    - \gamma \nabla Q_n(\Phi),
\]
where each summand is given by the symmetric operator associated
to the corresponding quadratic form. The operator
$L_n := \alpha \Delta_n^{\rm total} + \beta
\mathrm{d}_n^* \mathrm{d}_n - \gamma I$
(with $\Delta_n^{\rm total}$ the direct sum of $\Delta_n$ on
$X_n^0$ and $\partial_n^* \partial_n$ on $X_n^1$, and
$\mathrm{d}_n^* \mathrm{d}_n$ extended by zero on $X_n^1$) is
bounded and self-adjoint, and $\nabla F_n(\Phi) = L_n \Phi$. Both
$F_n$ and $\nabla F_n$ are $W$-equivariant and
$\sigma$-equivariant.
\end{lemma}

\begin{proof}
Immediate from Definition~\ref{def:Fn}: each quadratic form
$\Phi \mapsto \langle \Phi, A\Phi \rangle$ with $A$ symmetric is
continuously differentiable with gradient $2 A\Phi$, so
$\nabla F_n(\Phi) = L_n \Phi$ as stated. Self-adjointness of $L_n$
follows from self-adjointness of each constituent operator
(combinatorial Laplacian is self-adjoint by inspection;
$\mathrm{d}_n^*\mathrm{d}_n$ is self-adjoint by construction;
$I$ is trivially self-adjoint). $W$- and $\sigma$-equivariance of
$L_n$ follow from $W$- and $\sigma$-equivariance of the
underlying combinatorial structures (Lemma~\ref{lem:commute-subst}).
\end{proof}

\subsection{Finite-level cascade gradient flow}
\label{subsec:gradient-flow}

\begin{definition}[Finite-level cascade gradient flow]
\label{def:gradient-flow}
The \emph{continuous-time finite-level cascade gradient flow} on
$X_n^0 \oplus X_n^1$ is the ordinary differential equation
\begin{equation}
\label{eq:flow}
  \frac{d \Phi_n(t)}{dt} = -\nabla F_n(\Phi_n(t))
  = -L_n \Phi_n(t), \qquad t \geq 0.
\end{equation}
The \emph{discrete-time cascade gradient flow} with step size
$\eta_n > 0$ is the iteration
\begin{equation}
\label{eq:flow-discrete}
  \Phi_{n, k+1} = \Phi_{n, k} - \eta_n L_n \Phi_{n, k},
  \qquad k = 0, 1, 2, \ldots,
\end{equation}
starting from any $\Phi_{n, 0} \in X_n^0 \oplus X_n^1$.
\end{definition}

\begin{theorem}[Existence, uniqueness, and energy decrease]
\label{thm:flow-exists}
Fix $n \geq 0$ and an initial datum $\Phi_{n, 0} \in X_n^0 \oplus
X_n^1$.
\begin{enumerate}
\item The continuous-time flow \eqref{eq:flow} admits a unique
  global-in-time $C^1$ solution $\Phi_n \colon [0, \infty) \to X_n^0
  \oplus X_n^1$ with $\Phi_n(0) = \Phi_{n, 0}$, explicitly
  $\Phi_n(t) = e^{-t L_n} \Phi_{n, 0}$.
\item The closure energy is non-increasing along the flow:
  $F_n(\Phi_n(t_2)) \leq F_n(\Phi_n(t_1))$ for $t_1 \leq t_2$, and
  strictly decreasing unless $\Phi_n(t_1) \in \ker L_n$.
\item For any step size $\eta_n$ satisfying
  $0 < \eta_n \leq \|L_n\|^{-1}$ (operator norm), the discrete-time
  iteration \eqref{eq:flow-discrete} converges to the nearest
  critical point $\Phi_n^* \in \ker L_n$ in the $X_n^0 \oplus
  X_n^1$ norm, and $F_n(\Phi_{n, k})$ is non-increasing in $k$.
\item The flow is $W$-equivariant and $\sigma$-equivariant in the
  following sense: if $g$ is any $W$-element or $\sigma$, then
  $\Phi_n(0) \mapsto g\Phi_n(0)$ produces the flow
  $t \mapsto g\Phi_n(t)$.
\end{enumerate}
\end{theorem}

\begin{proof}
\emph{(1)} Since $L_n$ is a bounded linear operator on a
finite-dimensional Hilbert space, $-L_n$ generates a $C^0$-semigroup
$(e^{-tL_n})_{t \geq 0}$ of contractions (when
$\alpha, \gamma$ are chosen such that $L_n$ is bounded below; if
$L_n$ has negative eigenvalues the semigroup may grow exponentially,
but existence is unchanged). Uniqueness follows from the
Picard--Lindel\"of theorem for linear ODEs.

\emph{(2)} Along the flow,
\[
  \frac{d}{dt} F_n(\Phi_n(t))
  = \langle \nabla F_n(\Phi_n(t)), \dot\Phi_n(t) \rangle
  = \langle L_n \Phi_n(t), -L_n \Phi_n(t) \rangle
  = -\|L_n \Phi_n(t)\|^2 \leq 0,
\]
with equality iff $L_n \Phi_n(t) = 0$, i.e.\ $\Phi_n(t) \in \ker
L_n$. This gives non-increase and strict decrease off the kernel.

\emph{(3)} For $0 < \eta_n \leq \|L_n\|^{-1}$, the iteration map
$T_{\eta_n}(\Phi) := \Phi - \eta_n L_n \Phi$ has operator norm
$\|I - \eta_n L_n\| \leq 1$, with strict contraction in the
orthogonal complement of $\ker L_n$. Standard finite-dimensional
gradient-descent theory gives convergence to the orthogonal
projection of $\Phi_{n, 0}$ onto $\ker L_n$, and energy
non-increase follows from
\[
  F_n(\Phi_{n, k+1}) - F_n(\Phi_{n, k}) = -\eta_n \|L_n
    \Phi_{n, k}\|^2 + O(\eta_n^2),
\]
where the $O(\eta_n^2)$ is non-positive for $\eta_n \leq
\|L_n\|^{-1}$.

\emph{(4)} Since $L_n$ is $W$- and $\sigma$-equivariant
(Lemma~\ref{lem:Fn-regularity}), so is $e^{-t L_n}$, hence so is
the flow.
\end{proof}

\begin{remark}[Scope of the flow]
Theorem~\ref{thm:flow-exists} is a finite-dimensional
gradient-flow theorem; no analytic or geometric continuum limit is
invoked. The \emph{refinement} behaviour of the flow as $n \to
\infty$, and its relation to classical continuum PDE flows (Ricci
flow, Navier--Stokes, YM heat flow), is the subject of the
conditional theorems in Sections~\ref{sec:metric-hydro},
\ref{sec:gauge}, and \ref{sec:handoff}. At the present stage the
flow is well-defined but \emph{uninterpreted}.
\end{remark}

\subsection{Full symmetry actions on $X_\cas$}
\label{subsec:sigma-full}

\begin{definition}[$\sigma$-action on the full configuration space]
\label{def:sigma-X-cas}
The $\sigma$-action of Definition~\ref{def:sigma-action-subst}
extends pointwise from $X_n^0$ to $X_\cas$: for
$\Phi = (\Phi^0, \Phi^1) \in X_\cas$, define
\[
  \sigma \Phi := (\sigma \Phi^0, \sigma \Phi^1),
\]
where $\sigma$ acts trivially on $X_n^1$ (since $\Im(\OO)$ carries
no non-trivial $\Q(\pphi)$-structure that survives the trivial
extension from $\Q$). This defines a unitary involution
$\sigma \colon X_\cas \to X_\cas$.
\end{definition}

\begin{lemma}[Commutation, inverse-limit version]
\label{lem:commute}
On $X_\cas$, the Coxeter actions, $\sigma$, the refinement shift,
and the finite-level flow $\Phi \mapsto e^{-tL_n}\Phi$ (at any
level $n$) all commute pairwise on their common domains of
definition. The sector-restriction projections $P_{H_4}, P_{D_4}$
(orthogonal projection onto the $H_4$ or $D_4$ sector of
$V_{H_4}\oplus V_{D_4}\oplus V_{16}\oplus V_8$) also commute with
all three symmetry actions.
\end{lemma}

\begin{proof}
Commutation of Coxeter, $\sigma$, and refinement on $X_n^0$ and
$X_n^1$ is Lemma~\ref{lem:commute-subst}. The inverse-limit and
admissible-subspace constructions preserve commutation by
naturality. Commutation with the flow $e^{-tL_n}$ follows from
$W$- and $\sigma$-equivariance of $L_n$
(Lemma~\ref{lem:Fn-regularity}). Commutation of the
sector-restriction projections with $W$ and refinement is by
construction of Definition~\ref{def:fibres}; commutation with
$\sigma$ is because $\sigma$ respects the sector decomposition
(trivial on $V_{D_4}, V_{16}, V_8$; non-trivial only on $V_{H_4}$).
\end{proof}

\begin{remark}[Rung restriction]
The sector-restriction projections $P_{H_4}, P_{D_4}$ are the
precise Hilbert-space version of what
\cite[\S\S 2--3]{CascadeEmbeddings} calls ``$\pi_H$'' (projection
onto $H_4$-rung content) and ``$\iota_D$'' (inclusion of
$D_4$-rung content into the $H_4$-rung). We adopt the
$P_{H_4}, P_{D_4}$ notation throughout; the identifications are
$\pi_H \leftrightarrow P_{H_4}$ and $\iota_D \leftrightarrow
\iota_{V_{D_4} \hookrightarrow V_{H_4}}$, the latter being the
obvious coordinate inclusion.
\end{remark}

%=====================================================================
\section{Abstract Correspondence Scheme}
\label{sec:corresp}
%=====================================================================

This section isolates the common structure used by all seven
projections constructed in the following sections. A
\emph{correspondence} is defined abstractly as a measurable,
symmetry-intertwining map from an admissible subspace of a
cascade sector space into a classical target category. The main
technical result of this section,
Theorem~\ref{thm:abstract-correspondence}, is that a map
satisfying these axioms automatically inherits the finer
properties (Borel measurability on $X_\cas$, equivariance on the
symmetric kernel, compatibility with refinement) used in the
downstream sector theorems. The purpose of the abstraction is to
avoid repeating the same verification six times in
Sections~\ref{sec:metric-hydro}--\ref{sec:cohomology-comp}.

\subsection{Target categories}
\label{subsec:target-cats}

\begin{definition}[Classical target]
\label{def:target}
A \emph{classical target} is a triple
$\mathcal{T} = (T, \tau_T, S_T)$ where:
\begin{itemize}
\item $T$ is a topological space (the \emph{target space});
\item $\tau_T$ is a Borel $\sigma$-algebra on $T$;
\item $S_T$ is a Borel-measurable right action of a group $G_T$
  on $T$ (the \emph{target symmetry group}).
\end{itemize}
Examples: $\mathcal{T}_{\met} = (\Met(M), {\rm Borel}, {\rm
Diff}(M))$ for a fixed smooth manifold $M$; $\mathcal{T}_{\gge} =
(\Conn(P), {\rm Borel}, {\rm Gauge}(P))$ for a fixed principal
$G$-bundle $P$; $\mathcal{T}_\zeta = (\mathcal{B}(L^2(\R_+, dx/x)),
{\rm strong Borel}, \{1\})$ for the bounded-operator target of
$T_\zeta[h]$; and so on.
\end{definition}

\begin{definition}[Cascade source]
\label{def:source}
A \emph{cascade source} is a pair
$\mathcal{S} = (D, G_D)$ where $D \subseteq X_\cas$ is an
admissible subspace of the cascade configuration space (a Borel
subset closed under the symmetry actions of
Section~\ref{sec:dynamics}) and $G_D$ is the restriction to $D$ of
$W \times \langle \sigma \rangle \times {\rm Refine}$, the product
of the three commuting symmetry actions of
Lemma~\ref{lem:commute}.
\end{definition}

\begin{definition}[Correspondence]
\label{def:correspondence}
A \emph{correspondence}
$\Pi \colon \mathcal{S} \rightsquigarrow \mathcal{T}$ is a
Borel-measurable map $\Pi \colon D \to T$ together with a
group homomorphism $\rho_\Pi \colon G_D \to G_T$ such that
\[
  \Pi(g \cdot \Phi) = \rho_\Pi(g) \cdot \Pi(\Phi) \quad
    \text{for all } \Phi \in D,\, g \in G_D
\]
(symmetry intertwining).
\end{definition}

\begin{remark}[Interpretation of $\rho_\Pi$]
The homomorphism $\rho_\Pi$ translates cascade symmetries into
classical symmetries. For the metric projection $\Pi_{\met}$, the
$W(D_4)$-action on the cascade side maps into an isometry of the
target metric (rotations of $\R^4$); for the gauge projection
$\Pi_{\gge}$, the cascade symmetries map into gauge
transformations; for $\Pi_\sigma$, the $\sigma$-involution maps
into the Galois-conjugation symmetry of $\Q(\pphi)$-cohomology.
\end{remark}

\subsection{General theorem}
\label{subsec:general-thm}

\begin{theorem}[Abstract correspondence lemma]
\label{thm:abstract-correspondence}
Let $\Pi \colon \mathcal{S} \rightsquigarrow \mathcal{T}$ be a
correspondence as in Definition~\ref{def:correspondence}, and
suppose further that:
\begin{enumerate}
\item[(A1)] $\Pi$ is the limit of finite-level maps
  $\Pi^{(n)} \colon D \cap (X_n^0 \oplus X_n^1) \to T$ in the
  following sense: for each $\Phi \in D$ and each level $n$,
  $\Pi^{(n)}(\Phi_n)$ is defined, and
  $\Pi^{(n)}(\Phi_n) \to \Pi(\Phi)$ in $T$ as $n \to \infty$;
\item[(A2)] each $\Pi^{(n)}$ is continuous in the $X_n^0 \oplus
  X_n^1$ norm;
\item[(A3)] each $\Pi^{(n)}$ is symmetry-intertwining at its own
  level: $\Pi^{(n)}(g \cdot \Phi_n) = \rho_\Pi(g) \cdot
  \Pi^{(n)}(\Phi_n)$ for $g \in W \times \langle\sigma\rangle$;
\item[(A4)] the refinement map commutes with $\Pi^{(n)}$ up to the
  inclusion of targets: there is a compatible system of target maps
  $T^{(n)} \hookrightarrow T$ making the diagram commute.
\end{enumerate}
Then:
\begin{enumerate}
\item $\Pi$ is Borel-measurable on $D$;
\item $\Pi$ is continuous on $D$ with the inverse-limit topology;
\item $\Pi$ is fully symmetry-intertwining.
\end{enumerate}
\end{theorem}

\begin{proof}
(1) Borel measurability: each $\Pi^{(n)}$ is continuous
(hence Borel) on its finite-dimensional domain by (A2). The inverse
limit of Borel maps with compatible target inclusions is Borel on
the inverse limit, by the projective-limit topology theorem
(\cite[Thm.~1.8.9]{Bogachev}).

(2) Continuity: follows from (A2) and (A1): pointwise convergence
of Borel-measurable functions that are each continuous, combined
with the inverse-limit topology, gives continuity of the limit
function on the inverse-limit topology.

(3) Symmetry intertwining: (A3) plus (A4) give intertwining at
every level; passing to the limit gives intertwining on the
inverse limit. The action of the refinement shift commutes with
the correspondence by (A4).
\end{proof}

\begin{remark}[How this is used]
Theorem~\ref{thm:abstract-correspondence} is a technical workhorse:
each sector-specific projection in
Sections~\ref{sec:metric-hydro}--\ref{sec:cohomology-comp} is
constructed as a limit of finite-level maps, each satisfying (A1)--
(A4). Rather than reprove Borel measurability, continuity, and
equivariance sector by sector, we verify (A1)--(A4) sector by
sector and invoke
Theorem~\ref{thm:abstract-correspondence} to conclude.
\end{remark}

\begin{remark}[Conditional correspondences]
For the conditional correspondences of
Sections~\ref{sec:metric-hydro}--\ref{sec:cohomology-comp} (the
ones depending on $\mathbf{H_1}$--$\mathbf{H_4}$),
Theorem~\ref{thm:abstract-correspondence} still applies to the
unconditional \emph{finite-level} part $\Pi^{(n)}$; the conditional
hypothesis is needed only for the existence of the limit
$\Pi^{(n)}(\Phi_n) \to \Pi(\Phi)$ in a classical topology, which
is (A1). We therefore state each conditional theorem in the form
``assuming $\mathbf{H}_i$, (A1) holds, and hence (1)--(3) of
Theorem~\ref{thm:abstract-correspondence} hold''.
\end{remark}

%=====================================================================
\section{Metric and Hydrodynamic Correspondences}
\label{sec:metric-hydro}
%=====================================================================

This section constructs $\Pi_{\met}$ from $D_4$-sector cascade data
to metric data, and $\Pi_{\hyd}$ from $\sigma$-fixed $D_4$-sector
dynamics to velocity/pressure data. Both are built at finite level
from an explicit discrete moment tensor and shown to be
well-defined, continuous, and $W(D_4)$-equivariant. The continuum-
limit statements (smooth metric on a $3$-manifold; intertwining
with Navier--Stokes) are conditional on the geometric-limit
hypothesis $\mathbf{H_1}$ \emph{together with} several structural
hypotheses imported from the hardened downstream infrastructure
papers (P5, P6, P7) and one new foundations-level conjecture
$\mathbf{H_1'}$ controlling the extension of P4's mass-only
flow-intertwining to the full closure flow. These imports are
listed explicitly in \S\ref{subsec:metric-hydro-imports} before
any construction.

\subsection{Hypothesis imports from the hardened substrate (P4--P7)
  and the foundations-level conjecture $\mathbf{H_1'}$}
\label{subsec:metric-hydro-imports}

The following hypotheses are explicit inputs of every theorem in
\S\S~\ref{subsec:moment-tensor}--\ref{subsec:H1} below. Three are
restatements of hypotheses named in the hardened downstream
infrastructure papers; one is a new foundations-level conjecture
introduced here.

\begin{hypothesis}[Coord-injectivity of the refinement, imported from P5]
\label{hyp:coord-inj-import}
The coordinate map $V_n \to \R^4$, $v \mapsto x_v^{(n)}$, is
injective for every $n \geq 0$. \emph{Equivalent to} Hypothesis
``coordinate-injectivity of the refinement'' (Fact labelled as
hypothesis) of \cite[Hypothesis~fact:distinct]{CascadeMetricProjection},
which itself is verified at level $0$ and flagged as deferred-
for-all-$n$ in the status remark of that paper.
\end{hypothesis}

\begin{hypothesis}[Spherical refinement tiles $S^3$, imported from P6]
\label{hyp:tiling-import}
For every $n$, the spherical realization of the $D_4$-branch
refinement tiles $S^3$ (interior-disjoint cells whose union is
$S^3$). \emph{Equivalent to} \cite[Hypothesis~fact:tiling]{CascadeSchlafli},
verified at level $0$ and flagged as deferred-for-all-$n$ there.
Used in this section for the continuum-limit manifold
identification and for any GH-convergence statement.
\end{hypothesis}

\begin{hypothesis}[Uniform edge-length separation, imported from P6]
\label{hyp:separation-import}
There exists $c_e > 0$ such that $e_n \geq c_e h_n$ for all $n \geq
0$, where $e_n$ is the minimum edge length and $h_n$ the maximum
edge length of the level-$n$ $D_4$-refinement.
\emph{Equivalent to} \cite[Hypothesis~fact:separation]{CascadeSchlafli},
verified at level $0$ and flagged as deferred-for-all-$n$ there.
Used below for doubling / packing estimates in the continuum limit.
\end{hypothesis}

\begin{hypothesis}[Gradient-image compatibility for Leray decomposition, imported from P7]
\label{hyp:gradient-image-import}
Either (i) the image space $\mathcal V_n$ of the cascade-averaging
velocity map contains the image of the discrete gradient
$\nabla \mathcal Q_n$, so that $\mathcal V_n$ is closed under
Leray decomposition; or (ii) one works in the enlarged space
$\widetilde{\mathcal V}_n := \mathcal V_n + \nabla \mathcal Q_n$
throughout. \emph{Equivalent to}
\cite[Hypothesis~hyp:gradient-image]{CascadeHydroProjection}.
Used in any NS-target statement of $\Pi_{\hyd}$.
\end{hypothesis}

\begin{hypothesis}[$\mathbf{H_1'}$: full-flow extension of P4's mass-only intertwining]
\label{hyp:full-flow}
In the sense precisely controlled by \cite[\S 7]{CascadeClosureDynamics},
P4 proves refinement-intertwining of the cascade gradient operator
$L_n = \alpha A_n + \beta B_n - \gamma C_n$ only in the mass-only
regime $\alpha = \beta = 0$ (i.e.\ $L_n = -\gamma C_n$), and makes
no claim for arbitrary $(\alpha, \beta, \gamma)$. The foundations-
level conjecture $\mathbf{H_1'}$ asserts that this mass-only flow-
intertwining extends, in the continuum limit $n \to \infty$, to a
compatibility between the full discrete closure flow $\partial_t
\Phi_n = -L_n \Phi_n$ (arbitrary $\alpha, \beta, \gamma \geq 0$)
and the target continuum flow of
Theorems~\ref{thm:pimet-cond}--\ref{thm:pihyd-cond} below. $\mathbf{H_1'}$
is neither derived here nor in P4.
\end{hypothesis}

\begin{remark}[Status of $\mathbf{H_1'}$]
\label{rmk:Hprime-status}
$\mathbf{H_1'}$ is the honest statement of what must additionally
hold, beyond P4's mass-only result, for the Ricci-Laplacian-mass
PDE form (Theorem~\ref{thm:pimet-cond}) and the Navier-Stokes PDE
form (Theorem~\ref{thm:pihyd-cond}) to be the correct continuum
limits of the \emph{full} cascade gradient flow. Without $\mathbf{H_1'}$,
only the mass-only sub-case survives rigorously
(Remark~\ref{rmk:massonly-survives} below). Discharging $\mathbf{H_1'}$
requires establishing level-uniform refinement compatibility for
the $\alpha A_n$ and $\beta B_n$ blocks, a problem P4 explicitly
defers.
\end{remark}

\begin{remark}[Mass-only sub-case survives unconditionally on $\mathbf{H_1'}$]
\label{rmk:massonly-survives}
In the mass-only regime $\alpha = \beta = 0$, the flow-projection
portion of Theorems~\ref{thm:pimet-cond} and~\ref{thm:pihyd-cond}
reduces to the statement that the continuum flow is the mass-
rescaling $\partial_t g = 2\gamma g$ (or the analogous trivial
velocity rescaling for $\Pi_{\hyd}$). This does \emph{not} require
$\mathbf{H_1'}$: it follows directly from P4 Theorem
``refinement-intertwining in the mass-only case''
\cite[Thm.~thm:refinement-compat]{CascadeClosureDynamics}. The
Ricci, Laplacian, and NS nonlinear-and-viscous terms in the full-
flow theorems below are the portion that requires $\mathbf{H_1'}$.
\end{remark>

\subsection{The discrete moment tensor}
\label{subsec:moment-tensor}

\begin{definition}[Edge-weighted moment tensor]
\label{def:moment-tensor}
For $\Phi \in X_n^0 \oplus X_n^1$ and a real-valued edge weight
$W_\Phi \colon E^{\rm or}(G_n^{D_4}) \to \R$ depending measurably on
$\Phi$, define
\begin{equation}
\label{eq:moment-tensor}
  M_n(\Phi)(v)_{ij}
  := \sum_{w \sim v} (w_i - v_i)(w_j - v_j) \, W_\Phi(v, w),
\end{equation}
for $v \in V(G_n^{D_4})$ and $i, j \in \{1, 2, 3, 4\}$, using the
$24$-cell coordinates of \cite[\S 2]{CascadeEmbeddings}.
\end{definition}

\begin{definition}[Canonical edge weight]
\label{def:canonical-weight}
The \emph{canonical edge weight} is
\[
  W_\Phi(v, w)
  := \tfrac{1}{2} \big( \|P_{D_4} \Phi^0(v)\|^2_{V_{D_4}}
    + \|P_{D_4} \Phi^0(w)\|^2_{V_{D_4}} \big)
    + \|\Phi^1(v, w)\|^2_{\Im(\OO)}.
\]
This is non-negative, symmetric, $W(D_4)$-invariant, and continuous
in $\Phi$.
\end{definition}

\begin{proposition}[Point-source moment tensor]
\label{prop:moment-point-source}
Let $\Phi^{\rm pt}$ be the configuration
$\Phi^0(e_1) = e_1, \Phi^0(v) = 0$ for $v \neq e_1, \Phi^1 \equiv 0$.
Then $M_0(\Phi^{\rm pt})(v) \propto \hat r(v) \otimes \hat r(v)$ for
each $v$ adjacent to $e_1$, where $\hat r(v) = (v - e_1)/\|v - e_1\|$.
\end{proposition}

\begin{proof}
Direct substitution into \eqref{eq:moment-tensor}: only the edge
$(e_1, v)$ contributes, and
$M_0(\Phi^{\rm pt})(v)_{ij} = (v - e_1)_i (v - e_1)_j W_{\Phi^{\rm pt}}
(e_1, v)$, which is proportional to $\hat r(v) \otimes \hat r(v)$.
\end{proof}

\begin{remark}[Consistency with cascade-gr]
Proposition~\ref{prop:moment-point-source} reproduces the
theorem-grade formula of \cite[\S C4.1]{CascadeGR} from the explicit
moment-tensor formula used here.
\end{remark}

\subsection{The finite-level metric projection}
\label{subsec:pimet-n}

\begin{definition}[Piecewise-flat metric datum]
\label{def:pf-metric}
The space of \emph{piecewise-flat metric data} on $G_n^{D_4}$ is
\[
  \mathcal{M}_n := \Sym^+(\R^4)^{V(G_n^{D_4})},
\]
equipped with vertexwise convergence in the Frobenius norm.
\end{definition}

\begin{definition}[Finite-level metric projection]
\label{def:pimet-n}
Let $\{\epsilon_n\}_{n \geq 0}$ be a positive sequence converging
to $0$ slower than the refinement scale. Define
$\Pi_{\met}^{(n)} \colon X_n^0 \oplus X_n^1 \to \mathcal{M}_n$ by
\[
  \Pi_{\met}^{(n)}(\Phi)(v) := M_n(\Phi)(v) + \epsilon_n I.
\]
\end{definition}

\begin{theorem}[$\Pi_{\met}^{(n)}$: well-defined, continuous, equivariant]
\label{thm:pimet-n}
For each $n \geq 0$, $\Pi_{\met}^{(n)}$ is (i) well-defined with
image in $\mathcal{M}_n$, (ii) locally Lipschitz in
$X_n^0 \oplus X_n^1$, (iii) $W(D_4)$-equivariant with $\rho_{\met}(g)
\cdot M := g M g^T$, and (iv) $\sigma$-invariant.
\end{theorem}

\begin{proof}
(i) Positive-semidefiniteness:
$\xi^T M_n(\Phi)(v) \xi = \sum_w (\xi \cdot (w - v))^2 W_\Phi(v, w)
\geq 0$ since $W_\Phi \geq 0$; adding $\epsilon_n I$ gives strict
positivity.

(ii) For $\Phi, \Psi$ in a bounded ball of radius $R$, the
Frobenius-norm difference
$\|M_n(\Phi) - M_n(\Psi)\|_F$ is controlled by
$\sum_{v \sim w} \|w - v\|^2 \cdot C_R \|\Phi - \Psi\|$ for a
constant $C_R$ depending on $R$ (from the quadratic continuity of
$W_\Phi$); this gives Lipschitz continuity on bounded sets.

(iii) With $g \in W(D_4)$ acting by $(g\Phi^0)(v) = g\Phi^0(g^{-1}v)$,
reindexing the sum in \eqref{eq:moment-tensor} gives
$M_n(g\Phi)(v) = g M_n(\Phi)(g^{-1}v) g^T$.

(iv) $\sigma$ acts trivially on $V_{D_4}, V_{16}, V_8, \Im(\OO)$, so
$W_{\sigma\Phi} = W_\Phi$ and hence $M_n(\sigma\Phi) = M_n(\Phi)$.
\end{proof}

\subsection{The finite-level hydrodynamic projection}
\label{subsec:pihyd-n}

\begin{definition}[Averaging kernel]
\label{def:avg-kernel}
Fix $\ell_{\rm cg} > 0$ and a smooth compactly-supported kernel
$K \colon \R^4 \to \R_{\geq 0}$ with $\int K = 1$ and support in
$B_{\ell_{\rm cg}}$. For $f \colon V(G_n^{D_4}) \to \R^k$,
define the averaged field
\[
  (A_{\ell_{\rm cg}} f)(x)
  := \frac{\sum_v K(\ell_{\rm cg}^{-1}(x - v)) f(v)}
       {\sum_v K(\ell_{\rm cg}^{-1}(x - v))}
  \in \R^k.
\]
\end{definition}

\begin{definition}[Finite-level hydrodynamic projection]
\label{def:pihyd-n}
For a $C^1$-time-dependent $\Phi \colon [0, T] \to X_n^0 \oplus
X_n^1$, define $\Pi_{\hyd}^{(n)}(\Phi) = (u_n, p_n)$ by
\begin{align*}
  u_n(x, t) &:= A_{\ell_{\rm cg}} \big[\partial_t (P_{D_4}\Phi^0)
    \cdot e_{\rm rad}\big](x, t), \\
  p_n(x, t) &:= A_{\ell_{\rm cg}} \big[ \|P_{D_4} \Phi^0\|^2/2 \big](x, t),
\end{align*}
where $e_{\rm rad}(v) = v/\|v\|$.
\end{definition}

\begin{theorem}[$\Pi_{\hyd}^{(n)}$: basic properties]
\label{thm:pihyd-n}
$\Pi_{\hyd}^{(n)}$ is well-defined, continuous in the natural
$C^1 \to C^0$ norm, $W(D_4)$-equivariant (with the $O(4)$-action
on $(u, p)$), and $\sigma$-invariant.
\end{theorem}

\begin{proof}
The smoothness of $A_{\ell_{\rm cg}}$ gives well-definedness and
continuity. Equivariance follows from $W(D_4)$-equivariance of
$P_{D_4}$, the $O(4)$-invariance of $\|\cdot\|$, and the
translation-covariance of $A_{\ell_{\rm cg}}$. $\sigma$-invariance
is because $\sigma$ is trivial on $V_{D_4}$.
\end{proof}

\subsection{Inverse-limit metric projection (unconditional)}
\label{subsec:pimet-pihyd-inv}

\begin{proposition}[Inverse-limit metric projection, conditional on Hypothesis~\ref{hyp:coord-inj-import}]
\label{prop:pimet-inv-lim}
\emph{Assume Hypothesis~\ref{hyp:coord-inj-import}.} Under this
hypothesis, the maps $\Pi_{\met}^{(n)}$ satisfy (A2)--(A4) of
Theorem~\ref{thm:abstract-correspondence} for the $\sigma$- and
$W(D_4)$-actions and for refinement, and they form a compatible
family. They therefore induce a Borel measurable,
$\sigma$-invariant, $W(D_4)$-equivariant map
\[
  \Pi_{\met} \colon X_\cas^{\rm adm} \to \mathcal{M}_\infty,
  \qquad \mathcal{M}_\infty := \varprojlim_n \mathcal{M}_n.
\]
The proposition is \emph{not} unconditional: coordinate-
injectivity of the refinement is needed for the moment-tensor
kernel to be well-defined at all distinct vertex pairs (the
denominators $\|x_w - x_v\|$ entering the edge-weight and kernel
constructions are non-vanishing at distinct vertices only under
Hypothesis~\ref{hyp:coord-inj-import}). It is unconditional on the
continuum-limit hypothesis $\mathbf{H_1}$ and on the full-flow
conjecture $\mathbf{H_1'}$.
\end{proposition}

\begin{proof}
Under Hypothesis~\ref{hyp:coord-inj-import}: (A2) by
Theorem~\ref{thm:pimet-n}(ii); (A3) by
Theorem~\ref{thm:pimet-n}(iii)(iv); (A4) by the zero-extension
property of the refinement map $R_n$ (the moment tensor of the
zero-extended configuration is the zero-extension of the moment
tensor). Theorem~\ref{thm:abstract-correspondence} gives the three
conclusions.
\end{proof}

\subsection{Conditional correspondence under $\mathbf{H_1}$}
\label{subsec:H1}

\begin{hypothesis}[$\mathbf{H_1}$: Geometric limit]
\label{hyp:H1}
There exist (i) a smooth closed $3$-manifold $M$; (ii) a continuous
map $\pi_{\rm lim} \colon \mathcal{M}_\infty \to \mathcal{M}(M)$
sending inverse-limit piecewise-flat metric data to smooth
Riemannian metrics in the $C^\infty$ topology; (iii) an admissible
subspace $D_{\met} \subset X_\cas^{\rm adm}$ closed under $W(D_4)$,
$\sigma$, and refinement, on which $\pi_{\rm lim} \circ \Pi_{\met}$
is well-defined and continuous; (iv) verification of the
Burago--Ivanov and Cheeger--Colding hypotheses and a uniform bound
on the $g/\mathrm{geo}$-ratio overshoot flagged in
\cite[\S C2.bis]{CascadeGR} (items (a)--(d) of the open-item list
reproduced in \cite{MillenniumPoincare}).
\end{hypothesis}

\begin{remark}[Status of $\mathbf{H_1}$]
$\mathbf{H_1}$ is \emph{open}. Numerical evidence at depth
$n \leq 1$ and a proof sketch are in \cite[\S C2.bis]{CascadeGR};
four specific open items remain.
\end{remark}

\begin{theorem}[$\Pi_{\met}$ continuum correspondence, conditional
on $\mathbf{H_1}$ + full import chain + $\mathbf{H_1'}$ (full-flow portion)]
\label{thm:pimet-cond}
Assume Hypotheses~\ref{hyp:coord-inj-import}, \ref{hyp:tiling-import},
\ref{hyp:separation-import}, and $\mathbf{H_1}$. Then
$\Pi_{\met} := \pi_{\rm lim} \circ \Pi_{\met}^{\rm inv-lim}
\colon D_{\met} \to \mathcal{M}(M)$ is continuous, Borel,
$W(D_4)$-equivariant, and $\sigma$-invariant (\emph{the
well-definedness / continuity / equivariance portion; this part
does not require $\mathbf{H_1'}$}).

\emph{The flow-projection portion splits into two cases}:
\begin{enumerate}
\item[\textup{(a)}] \emph{Mass-only sub-case} $\alpha = \beta = 0$
  (no extra hypothesis beyond the well-definedness chain above):
  the cascade gradient flow projects, as $n \to \infty$, to the
  continuum mass-rescaling $\partial_t g = 2\gamma g$, by direct
  application of \cite[Thm.~thm:refinement-compat]{CascadeClosureDynamics}
  (P4's mass-only flow-intertwining) followed by continuum limit
  under $\mathbf{H_1}$.
\item[\textup{(b)}] \emph{Full-flow case} (arbitrary $\alpha,
  \beta, \gamma \geq 0$), \emph{conditional additionally on
  $\mathbf{H_1'}$} (Hypothesis~\ref{hyp:full-flow}): the cascade
  gradient flow projects to a continuum flow of the schematic
  form $\partial_t g = -2(\alpha \Ric + \beta \nabla^*\nabla -
  \gamma g)$ up to $O(\epsilon_n)$ corrections.
\end{enumerate}
The Ricci and Laplacian terms in (b) are conditional on
$\mathbf{H_1'}$ beyond what P4 supplies; without $\mathbf{H_1'}$
the full-flow PDE form is not established.
\end{theorem}

\begin{proof}
\emph{Well-definedness / continuity / equivariance.}
Under Hypotheses~\ref{hyp:coord-inj-import},
\ref{hyp:tiling-import}, \ref{hyp:separation-import}, and
$\mathbf{H_1}$(i)--(iv), hypothesis (A1) of
Theorem~\ref{thm:abstract-correspondence} is satisfied for the
family $\Pi_{\met}^{(n)}$ with target inclusion
$\mathcal{M}_n \hookrightarrow \mathcal{M}(M)$ via $\pi_{\rm lim}$.
Combining with Proposition~\ref{prop:pimet-inv-lim}, the abstract
correspondence lemma gives continuity, Borel measurability, and
equivariance.

\emph{Mass-only sub-case (a).} With $\alpha = \beta = 0$, the
discrete gradient operator is $L_n = -\gamma C_n$. By P4
Theorem~thm:refinement-compat \cite{CascadeClosureDynamics},
$p_{n+1, n} L_{n+1} = L_n p_{n+1, n}$, so the flow passes through
refinement. Under $\mathbf{H_1}$'s continuum limit, this limits to
$\partial_t g = 2\gamma g$.

\emph{Full-flow case (b), conditional on $\mathbf{H_1'}$.} Under
$\mathbf{H_1'}$ the mass-only intertwining extends to the full
closure flow in the $n \to \infty$ limit. Differentiating
\eqref{eq:moment-tensor} under the full flow
\eqref{eq:flow} and passing through $\pi_{\rm lim}$, the second-
order differences in $v$ limit to $\Ric$ (requires
$\mathbf{H_1'}$ for the $\alpha A_n$-block continuum extension),
the first-order differences to $\nabla^* \nabla$ (requires
$\mathbf{H_1'}$ for the $\beta B_n$-block extension), and the
mass-term $\gamma Q_n$ to $\gamma g$ (from the mass-only sub-case).
The $\epsilon_n I$ term contributes an $O(\epsilon_n)$ correction.
\end{proof}

\begin{theorem}[$\Pi_{\hyd}$ continuum correspondence, split into well-definedness, mass-only, and full-flow portions]
\label{thm:pihyd-cond}
Assume Hypotheses~\ref{hyp:coord-inj-import},
\ref{hyp:tiling-import}, \ref{hyp:separation-import},
\ref{hyp:gradient-image-import}, and $\mathbf{H_1}$. On an
admissible subspace $D_{\hyd} \subset X_\cas^{\rm adm}$ of
$\sigma$-fixed, $C^1$ time-dependent configurations, and with
P7's macroscopic-time-scaling convention $\tau = \lambda_n t$
\cite[Thm.~thm:projected-dynamics]{CascadeHydroProjection}:

\begin{enumerate}
\item[\textup{(i)}] \emph{Well-definedness / continuity /
  equivariance} (no $\mathbf{H_1'}$ required). The limit
  $\Pi_{\hyd} \colon D_{\hyd} \to C^0(\R_+ \times M; TM \oplus \R)$
  is well-defined, continuous, Borel, and $W(D_4)$-equivariant.
\item[\textup{(ii)}] \emph{Mass-only sub-case}
  $\alpha = \beta = 0$ (no $\mathbf{H_1'}$ required). The
  projected flow is the trivial rescaling $\partial_\tau u =
  2\gamma u$, $\partial_\tau p = 2\gamma p$, by direct application
  of \cite[Thm.~thm:refinement-compat]{CascadeClosureDynamics}
  (P4's mass-only flow-intertwining) followed by continuum limit.
\item[\textup{(iii)}] \emph{Full-flow case} (arbitrary
  $\alpha, \beta, \gamma \geq 0$), \emph{additionally conditional
  on $\mathbf{H_1'}$} (Hypothesis~\ref{hyp:full-flow}). The cascade
  gradient flow restricted to $D_{\hyd}$ projects to the
  macroscopic Navier--Stokes system on $M$ up to controlled error
  terms.
\end{enumerate}
Without $\mathbf{H_1'}$, only (i)--(ii) survive; the full-flow
NS PDE projection in (iii) is the portion that requires
$\mathbf{H_1'}$ beyond P4.
\end{theorem}

\begin{proof}[Proof (sketch)]
\emph{Well-definedness / continuity.} $\mathbf{H_1}$ supplies the
manifold $M$; the averaging kernel $A_{\ell_{\rm cg}}$ gives $C^0$
fields. Hypothesis~\ref{hyp:gradient-image-import} ensures the
Leray projector $\Pi_n$ of P7 is well-defined on the velocity
image space. Hypothesis~\ref{hyp:coord-inj-import} ensures the
moment-tensor denominators used in the averaging are non-vanishing;
Hypotheses~\ref{hyp:tiling-import}, \ref{hyp:separation-import}
supply the GH-convergence and doubling estimates of the continuum
limit.

\emph{Flow projection to NS, under $\mathbf{H_1'}$.} The
identification of $\partial_\tau(P_{D_4}\Phi^0)$ with velocity at
macroscopic time requires the $\lambda_n^{-1}$ factor of
\cite[Thm.~thm:projected-dynamics, $D_n^{\rm comm}$]{CascadeHydroProjection}.
Under $\mathbf{H_1'}$, the full closure flow
$\partial_t \Phi_n = -L_n \Phi_n$ (arbitrary $\alpha, \beta,
\gamma$) extends P4's mass-only intertwining, and standard coarse-
graining identities relating discrete gradient/divergence to
continuum analogues give: the nonlinear $(u \cdot \nabla)u$ and
pressure gradient terms from the quadratic structure of $F_n$, and
the viscous $\nu\Delta u$ term from $\alpha R_n$. Explicit PDE
estimates of the convergence rate as $n\to\infty$ and
$\ell_{\rm cg}\to 0$ are part of the open items of $\mathbf{H_1}$
and $\mathbf{H_1'}$ and not discharged here.

\emph{Mass-only sub-case.} Under $\alpha = \beta = 0$, P4's
refinement-intertwining is unconditional (no $\mathbf{H_1'}$); the
continuum limit under $\mathbf{H_1}$ gives the trivial rescaling
statement.
\end{proof}

\begin{remark}[Scope limits]
Theorem~\ref{thm:pihyd-cond} does not resolve Navier--Stokes. It
gives a structural correspondence, with the full-flow PDE
projection conditional on the full hypothesis chain
Hypotheses~\ref{hyp:coord-inj-import}--\ref{hyp:gradient-image-import}
+ $\mathbf{H_1}$ + $\mathbf{H_1'}$. Only the mass-only sub-case
survives without $\mathbf{H_1'}$. The downstream $L^\infty$
vorticity bound and BKM application in \cite{MillenniumNS} require
the open items of $\mathbf{H_1}$, $\mathbf{H_1'}$, \emph{and} the
P5/P6/P7 structural hypotheses, plus additional PDE machinery.
\end{remark}

%=====================================================================
\section{Gauge Correspondence}
\label{sec:gauge}
%=====================================================================

This section constructs the gauge projection $\Pi_{\gge}$ from
octonionic edge data (the edge-sector $X_n^1$ with $\Im(\OO)$
fibre, Definition~\ref{def:X-n-1}) to cylindrical $\mathrm{SU}(3)$
holonomy data. The finite-level map is a standard lattice-gauge
construction: exponentiate a Lie-algebra-valued $1$-cochain to get
a group-valued holonomy. The cascade-specific feature is that the
Lie algebra is $\mathrm{su}(3) \subset G_2 = \Der(\OO)$, selected
by the stabilizer of a fixed imaginary unit
\cite[\S 6]{CascadeEmbeddings}. The continuum-limit statement
(classical $L^2_1$ connection; YM action transport) is conditional
on the gauge-reconstruction hypothesis $\mathbf{H_2}$.

\subsection{Scope of the gauge construction: modeling choices}
\label{subsec:gauge-modeling}

Before any construction, we record three modeling choices that
together define what ``gauge action on cascade edge data'' means
in this section. Each is a choice of convention, not a derived
consequence of the classical $\mathrm{SU}(3) \subset G_2 \subset
\Der(\OO)$ package:

\begin{enumerate}
\item[\textup{(M1)}] \emph{Choice of fixed imaginary unit $i_0 \in
  \Im(\OO)$ with $i_0^2 = -1$.} The stabilizer subgroup of $i_0$
  under the $G_2$-action is a specific copy of $\mathrm{SU}(3)$ in
  $G_2$; different $i_0$ give different (but conjugate) copies.
  Different choices of $i_0$ produce different projections
  $P_{\mathfrak{su}(3)}$ and hence different
  $\mathfrak{su}(3)$-valued edge fields. We fix one $i_0$ once
  and for all; the downstream gauge quotient absorbs the
  conjugation ambiguity.

\item[\textup{(M2)}] \emph{Choice of injection
  $\iota_{\mathfrak{g}_2} \colon \Im(\OO) \to \mathfrak{g}_2 =
  \Der(\OO)$.} As noted in Definition~\ref{def:cascade-gauge-data},
  octonion commutators $[X, \cdot]_{\OO}$ are \emph{not} in
  general derivations of $\OO$; a non-trivial lift is needed to
  embed $\Im(\OO)$ into $\Der(\OO)$. The lift used here is the
  standard associator-respecting-derivation construction of
  \cite[\S 2]{BaezOctonions}; other lifts exist and give related
  but not identical projections. This is a modeling choice.

\item[\textup{(M3)}] \emph{Choice of the gauge action on
  octonionic edge data.} Classical lattice gauge theory takes
  edge data directly in the gauge group (here $\mathrm{SU}(3)$).
  Cascade edge data lives instead in $\Im(\OO)$ and must be
  connected to the gauge group via the
  lift-project-exponentiate chain $\Im(\OO) \xrightarrow{\iota_{\mathfrak{g}_2}}
  \mathfrak{g}_2 \xrightarrow{P_{\mathfrak{su}(3)}} \mathfrak{su}(3)
  \xrightarrow{\exp} \mathrm{SU}(3)$. The induced ``gauge
  transformation of $A \in \Im(\OO)$ by $u \in \mathrm{SU}(3)$''
  is therefore a paper-specific definition: it is whatever
  transformation of the $\mathfrak{g}_2$-lift $\iota_{\mathfrak{g}_2}(A)$
  agrees with the standard $\mathrm{SU}(3)$ adjoint action after
  $P_{\mathfrak{su}(3)}$-projection.
\end{enumerate}

All gauge-equivariance statements in
\S\S~\ref{subsec:pigauge-n}--\ref{subsec:H2} below are understood
as statements relative to these three modeling choices. A reader
preferring different choices would obtain a conjugate (but not
identical) construction.

\subsection{Octonionic gauge data}
\label{subsec:octo-gauge}

\begin{definition}[Fixed imaginary unit]
\label{def:i-fixed}
Fix a unit imaginary octonion $i_0 \in \Im(\OO)$ with
$i_0^2 = -1$. The stabilizer of $i_0$ under the $G_2 = \Der(\OO)$
action is the subgroup $\mathrm{SU}(3) \subset G_2$ acting on the
orthogonal complement $i_0^\perp \subset \Im(\OO)$, which carries a
natural complex structure via left-multiplication by $i_0$
\cite[\S 6]{CascadeEmbeddings}, \cite[\S 1]{BaezOctonions}.
Different choices of $i_0$ give conjugate copies of $\mathrm{SU}(3)$;
we fix one.
\end{definition}

\begin{notation}[Lie-algebra decomposition]
\label{not:g2-decomp}
Write $\mathfrak{g}_2 = \Der(\OO)$ for the $14$-dimensional Lie
algebra of $G_2$. The $\mathrm{SU}(3)$-stabilizer of
Definition~\ref{def:i-fixed} decomposes $\mathfrak{g}_2$ as
\[
  \mathfrak{g}_2 = \mathfrak{su}(3) \oplus \mathfrak{m},
\]
where $\mathfrak{su}(3)$ is the stabilizer algebra (real dimension
$8$) and $\mathfrak{m}$ is a $6$-dimensional complement stable
under the adjoint action of $\mathrm{SU}(3)$
\cite[\S 2.4]{BaezOctonions}. Let $P_{\mathfrak{su}(3)}$ denote the
orthogonal projection $\mathfrak{g}_2 \to \mathfrak{su}(3)$ with
respect to the Killing form.
\end{notation}

\begin{definition}[Cascade gauge datum]
\label{def:cascade-gauge-data}
A \emph{cascade gauge datum} at level $n$ is an element
$A \in X_n^1$ (cf.\ Definition~\ref{def:X-n-1}). Given such $A$,
the corresponding \emph{$\mathfrak{su}(3)$-valued edge field} is
\[
  A^{\mathfrak{su}(3)}(e)
  := P_{\mathfrak{su}(3)} \big( \iota_{\mathfrak{g}_2}(A(e)) \big)
  \in \mathfrak{su}(3),
  \qquad e \in E^{\rm or}(G_n^{\bullet}),
\]
where $\iota_{\mathfrak{g}_2} \colon \Im(\OO) \to \mathfrak{g}_2$
is the standard injection of imaginary octonions into the
$14$-dimensional Lie algebra $\mathfrak{g}_2 = \Der(\OO)$, and
$P_{\mathfrak{su}(3)}$ is the projection of
Notation~\ref{not:g2-decomp}. \emph{Note: the map
$\iota_{\mathfrak{g}_2}$ is not the naive octonion-commutator
$[A(e), \cdot]_{\OO}$ --- octonion commutators do not lie in
$\Der(\OO)$ in general, since the octonions are nonassociative.
The correct injection $\Im(\OO) \hookrightarrow \mathfrak{g}_2$
goes via the associator-respecting derivation lift, e.g.\ the
map $X \mapsto 3[L_X, L_Y] + ...$ formalism of
\cite[\S 2]{BaezOctonions}; we use this injection here without
reviewing its explicit form.}
\end{definition}

\begin{remark}[Construction, not a theorem of \cite{CascadeEmbeddings}]
Definition~\ref{def:cascade-gauge-data} is a construction of
this paper using the classical $\mathrm{SU}(3) \subset G_2$
stabilizer-of-$i_0$ picture of
\cite[\S 6]{CascadeEmbeddings} together with the injection
$\iota_{\mathfrak{g}_2}$ of associator-respecting derivations
\cite[\S 2]{BaezOctonions}. It is \emph{not} itself a theorem of
\cite{CascadeEmbeddings}: the precise algebraic projection
$P_{\mathfrak{su}(3)} \circ \iota_{\mathfrak{g}_2}$ is an
interpretive choice made here. Different $i_0$ give unitarily
equivalent projections.
\end{remark>

\subsection{The finite-level gauge projection}
\label{subsec:pigauge-n}

\begin{definition}[Cylindrical holonomy data]
\label{def:holonomy-n}
The space of \emph{cylindrical $\mathrm{SU}(3)$ holonomy data} on
$G_n^{\bullet}$ is
\[
  \mathcal{A}_n
  := \big\{ U \colon E^{\rm or}(G_n^{\bullet}) \to \mathrm{SU}(3)
     : U(-e) = U(e)^{-1} \big\},
\]
equipped with the product topology from the Lie-group topology on
$\mathrm{SU}(3)$.
\end{definition}

\begin{definition}[Finite-level gauge projection]
\label{def:pigauge-n}
Fix an edge length scale $a_n > 0$ with $a_n \to 0$ at the
refinement rate. The \emph{finite-level gauge projection}
$\Pi_{\gge}^{(n)} \colon X_n^1 \to \mathcal{A}_n$ is
\[
  \Pi_{\gge}^{(n)}(A)(e)
  := \exp\big( a_n \cdot A^{\mathfrak{su}(3)}(e) \big)
  \in \mathrm{SU}(3),
  \qquad e \in E^{\rm or}(G_n^{\bullet}).
\]
Here $\exp \colon \mathfrak{su}(3) \to \mathrm{SU}(3)$ is the
Lie-algebra exponential, and the reversal convention $A(-e) = -A(e)$
(Definition~\ref{def:X-n-1}) gives
$\Pi_{\gge}^{(n)}(A)(-e) = \exp(-a_n A^{\mathfrak{su}(3)}(e)) =
\Pi_{\gge}^{(n)}(A)(e)^{-1}$, as required for holonomy data.
\end{definition}

\begin{theorem}[$\Pi_{\gge}^{(n)}$: well-defined, continuous,
equivariant]
\label{thm:pigauge-n}
For each $n \geq 0$:
\begin{enumerate}
\item $\Pi_{\gge}^{(n)}$ is well-defined: for every $A \in X_n^1$,
  $\Pi_{\gge}^{(n)}(A) \in \mathcal{A}_n$.
\item $\Pi_{\gge}^{(n)}$ is continuous from the $L^2$ topology on
  $X_n^1$ to the product topology on $\mathcal{A}_n$, and is
  locally Lipschitz on bounded sets.
\item $\Pi_{\gge}^{(n)}$ is compatible with the $\mathrm{SU}(3)
  \subset G_2$ identification of
  \cite[\S 6]{CascadeEmbeddings}: the image of
  $\Pi_{\gge}^{(n)}(A)$ lies in the $\mathrm{SU}(3)$ subgroup
  selected by the stabilizer of $i_0$.
\item \emph{Gauge-equivariance, conditional on modeling choices
  (M1)--(M3) of \S\ref{subsec:gauge-modeling}.}
  Define the induced gauge action on cascade edge data $A \in
  \Im(\OO)$ by the lift-adjoint-project chain of (M3): for a
  lattice gauge transformation $u \colon V(G_n^{\bullet}) \to
  \mathrm{SU}(3)$,
  \[
    (u \cdot A)(e) \;:=\; P_{\mathfrak{su}(3)}\bigl(
      \mathrm{Ad}_{u(e_{\rm start})} \bigl[\iota_{\mathfrak{g}_2}(A(e))\bigr]
      \cdot u(e_{\rm end})^{-1}\bigr)^{-1}\!
    \text{ (read back via (M2) to } \Im(\OO) \text{)},
  \]
  interpreted to leading order in $a_n$ via BCH. Under this
  modeling-derived gauge action, the projection transforms as
  $\Pi_{\gge}^{(n)}(u \cdot A)(e) = u(e_{\rm start}) \cdot
  \Pi_{\gge}^{(n)}(A)(e) \cdot u(e_{\rm end})^{-1}$ up to $O(a_n^2)$
  corrections. The $O(a_n^2)$ corrections are the
  parallel-transport-curvature terms, and their vanishing is
  \emph{not} claimed here; they become the per-level defect
  absorbed in the inverse-limit hypothesis of
  Proposition~\ref{prop:pigauge-inv-lim}.
\end{enumerate}
\end{theorem}

\begin{proof}
(1) $A^{\mathfrak{su}(3)}(e) \in \mathfrak{su}(3)$ by
Definition~\ref{def:cascade-gauge-data}; the exponential of an
$\mathfrak{su}(3)$ element lies in $\mathrm{SU}(3)$. The reversal
convention is verified by direct computation using
$\exp(-X) = \exp(X)^{-1}$.

(2) Continuity: the projection $P_{\mathfrak{su}(3)}$ is
continuous (bounded linear) and $\mathrm{ad}$ is continuous, so
$A \mapsto A^{\mathfrak{su}(3)}$ is continuous; $\exp$ is smooth;
composition gives continuity. Local Lipschitz continuity on bounded
sets follows from the smoothness of $\exp$ on a compact subset of
$\mathfrak{su}(3)$.

(3) The $\mathrm{SU}(3)$ identification is by construction: the
stabilizer of $i_0$ is exactly the $\mathrm{SU}(3)$ subgroup whose
Lie algebra is the image of $P_{\mathfrak{su}(3)}$.

(4) Gauge equivariance is the standard lattice-gauge transformation
law, verified to leading order in $a_n$ by the
Baker--Campbell--Hausdorff formula. The $O(a_n^2)$ corrections are
the parallel-transport-curvature terms.
\end{proof}

\subsection{Inverse-limit gauge projection (unconditional)}
\label{subsec:pigauge-inv}

\begin{proposition}[Inverse-limit gauge projection, conditional on uniform lattice-curvature control]
\label{prop:pigauge-inv-lim}
Assume (as an additional lattice-gauge hypothesis) that the
$O(a_n^2)$ gauge-equivariance corrections of
Theorem~\ref{thm:pigauge-n}(4) admit uniform control across
refinement levels, so that the leading-order equivariance passes
to the inverse limit. Then the family
$\{\Pi_{\gge}^{(n)}\}_{n \geq 0}$ satisfies (A2)--(A4)
of Theorem~\ref{thm:abstract-correspondence} (for gauge-equivariance
replacing general symmetry equivariance), forms a compatible family
under refinement, and induces a Borel measurable, gauge-equivariant
map
\[
  \Pi_{\gge} \colon X_\cas^{\rm adm} \cap X^1_{\varprojlim}
  \to \mathcal{A}_\infty
  := \varprojlim_n \mathcal{A}_n.
\]
Without the uniform-curvature-control hypothesis, the conclusion
only holds at finite level with $O(a_n^2)$ residual corrections
at each level; the inverse-limit gauge-equivariance then requires
independent justification.
\end{proposition}

\begin{proof}
Under the uniform lattice-curvature-control hypothesis: (A2) by
Theorem~\ref{thm:pigauge-n}(2); (A3) by
Theorem~\ref{thm:pigauge-n}(3)(4) together with the assumed
uniform control, which lifts the level-$n$ $O(a_n^2)$
equivariance to exact equivariance in the limit. For (A4), the
refinement of a gauge datum is standard in lattice-gauge theory:
zero-extension of $A$ at the new edges gives $\Pi_{\gge}^{(n+1)}$
that restricts to $\Pi_{\gge}^{(n)}$ on the level-$n$ edges. The
abstract correspondence theorem gives the conclusions.
\end{proof>

\subsection{Conditional continuum-gauge correspondence under
$\mathbf{H_2}$}
\label{subsec:H2}

Let $\Conn(P)$ denote the space of $L^2_1$ (Sobolev-class) classical
$\mathrm{SU}(3)$-connections on a fixed principal
$\mathrm{SU}(3)$-bundle $P$ over a fixed smooth manifold $M$;
equip $\Conn(P)$ with the standard $L^2_1$ topology.

\begin{hypothesis}[$\mathbf{H_2}$: Gauge reconstruction]
\label{hyp:H2}
There exist:
\begin{enumerate}
\item a smooth manifold $M$ and a principal $\mathrm{SU}(3)$-bundle
  $P \to M$;
\item a continuous map
  $\pi_{\gge, {\rm lim}} \colon \mathcal{A}_\infty \to \Conn(P)$
  sending inverse-limit holonomy data to $L^2_1$ connections;
\item an admissible subspace
  $D_{\gge} \subset X_\cas^{\rm adm} \cap X^1_{\varprojlim}$ on
  which $\pi_{\gge, {\rm lim}} \circ \Pi_{\gge}$ is well-defined,
  continuous, gauge-equivariant, and satisfies the
  Osterwalder--Schrader positivity axiom for the induced YM measure
  \cite{OSBook};
\item Wilson-loop area-law estimates under refinement: the
  lattice-level Wilson loop expectation values converge to continuum
  area-law expectations as $n \to \infty$, as required for
  confinement.
\end{enumerate}
\end{hypothesis}

\begin{remark}[Status of $\mathbf{H_2}$]
$\mathbf{H_2}$ is \emph{open} and is the content of ``constructive
quantum YM existence'' in the Clay statement. Parts (i)--(ii) are
the continuum completion (standard in lattice-gauge theory but not
yet theorem-grade for $\mathrm{SU}(3)$); parts (iii)--(iv) are the
OS positivity and area-law steps flagged in
\cite[\S MP1.2]{CascadeMillennium} as among the open technical
items for YM.
\end{remark}

\begin{theorem}[$\Pi_{\gge}$ continuum correspondence, conditional
on $\mathbf{H_2}$]
\label{thm:pigauge-cond}
Assume $\mathbf{H_2}$. Then
$\Pi_{\gge} := \pi_{\gge, {\rm lim}} \circ \Pi_{\gge}^{\rm inv-lim}
\colon D_{\gge} \to \Conn(P)$ is continuous, Borel,
gauge-equivariant, and compatible with the
$\mathrm{SU}(3) \subset G_2$ identification. Moreover, the
restriction of the cascade closure functional $\mathcal{F}$ to the
gauge sector transports, under $\Pi_{\gge}$, to a constant
multiple of the classical Yang--Mills action
\[
  S_{YM}(A) = \tfrac{1}{2} \int_M \operatorname{tr}(F(A) \wedge *F(A)),
\]
with the proportionality constant determined by the
$\alpha, \beta, \gamma$ coefficients and the lattice-scale factor
$a_n$.
\end{theorem}

\begin{proof}[Proof (sketch)]
Under $\mathbf{H_2}$, (A1) of
Theorem~\ref{thm:abstract-correspondence} holds; combined with
Proposition~\ref{prop:pigauge-inv-lim}, the abstract lemma gives
continuity, measurability, and equivariance.

For the action transport, the cascade closure functional
restricted to gauge-sector data is, up to a quadratic form on
$\mathfrak{su}(3)$-valued edge fields, the discrete Yang--Mills
action in the standard lattice-gauge formulation
\cite{WilsonLatticeGauge}. In the $a_n \to 0$ continuum limit
(part of $\mathbf{H_2}$), the discrete YM action converges to the
continuous YM action $S_{YM}$; the proportionality constant comes
from the rescaling.

A Clay-level completion of this proof would require the full
Osterwalder--Schrader reconstruction and the Wilson-area-law
estimate, which are parts (iii)--(iv) of $\mathbf{H_2}$ and not
discharged here.
\end{proof}

\begin{remark}[Scope]
Theorem~\ref{thm:pigauge-cond} does not resolve the YM mass-gap
problem. It provides a structural correspondence from cascade
gauge data to classical $\mathrm{SU}(3)$ connections, with the
action transporting to the continuum YM action. The mass-gap bound
claimed in \cite{MillenniumYM} would still require $\mathbf{H_2}$
plus non-trivial spectral analysis.
\end{remark}

%=====================================================================
\section{Spectral--Arithmetic Correspondences}
\label{sec:spectral}
%=====================================================================

This section constructs the regularized transfer operators
$T_\zeta[h]$ (for the Riemann Hypothesis) and $T_E[h]$ (for
Birch--Swinnerton-Dyer). Both are built with an explicit cutoff
function $h$ of compact support; the cutoff removal $h \to
\mathbf{1}$ is the subject of the spectral-arithmetic trace
hypothesis $\mathbf{H_3}$. The finite-cutoff operators are
unconditionally well-defined as densely-defined symmetric or
bounded operators on standard Hilbert spaces.

The critical methodological point of this section: the earlier
Millennium RH/BSD manuscripts invoked raw operators $T_\zeta, T_E$
whose convergence is not established (Table~\ref{tab:exclusions}).
We replace those with regularized $T_\zeta[h], T_E[h]$, which are
well-defined, and state the cutoff-removal $h \to \mathbf{1}$ as
an explicit open hypothesis.

\subsection{The regularized spectral-zeta transfer operator}
\label{subsec:Tzeta-h}

\begin{definition}[Base operator $T_0$ on $L^2(\R_+, dx/x)$]
\label{def:T0}
Let $L^2(\R_+, dx/x)$ denote the Hilbert space of square-integrable
functions on $\R_+$ with multiplicative Haar measure
$dx/x$. Define
\[
  T_0 f(x) := \sum_{p \text{ prime}} \frac{\log p}{p^{1/2}}
    f(p \cdot x),
\]
with the sum interpreted as an abstract formal expression until
regularization is applied. At this stage $T_0$ is not a bounded
operator on $L^2(\R_+, dx/x)$; we regularize below.
\end{definition}

\begin{definition}[Prime-cutoff potential]
\label{def:Vsigma-h}
Let $h \colon \R \to \R$ be a real-valued, even, compactly-supported
smooth function. The \emph{prime-cutoff potential} is
\[
  V_{\sigma, h} f(x)
  := \sum_{p \text{ prime}} h(\log p) \cdot
    \frac{\log p}{p^{1/2}} \cdot f(p \cdot x),
\]
for $f \in C_c^\infty(\R_+)$. By the compact support of $h$, the
sum is finite: only finitely many primes contribute.
\end{definition}

\begin{definition}[Regularized transfer operator $T_\zeta[h]$, $\sigma$-symmetrized]
\label{def:Tzeta-h}
The \emph{regularized spectral-zeta transfer operator} is the
$\sigma$-symmetrization
\[
  T_\zeta[h] := \tfrac{1}{2}\bigl((T_0 \cdot \chi_h + V_{\sigma, h})
    + \sigma (T_0 \cdot \chi_h + V_{\sigma, h}) \sigma^{-1}\bigr),
\]
where $\sigma \colon L^2(\R_+, dx/x) \to L^2(\R_+, dx/x)$ is the
inversion $(\sigma f)(x) := f(1/x)$ (a unitary involution), and
$\chi_h$ is the multiplication operator by the indicator function
of $\{x \in \R_+ : \log x \in \operatorname{supp}(h)\}$. The
$\sigma$-symmetrization is explicit and built into the definition;
without it, the bare operator $T_0 \cdot \chi_h + V_{\sigma, h}$ is
\emph{not} symmetric on $L^2(\R_+, dx/x)$ (since $f \mapsto f(p x)$
has adjoint $f \mapsto f(x/p)$, not itself).
\end{definition}

\begin{theorem}[$T_\zeta[h]$: symmetric and bounded]
\label{thm:tzeta-h}
For $h$ as in Definition~\ref{def:Vsigma-h}, the operator
$T_\zeta[h]$ of Definition~\ref{def:Tzeta-h} is:
\begin{enumerate}
\item densely defined on $L^2(\R_+, dx/x)$, with core
  $C_c^\infty(\R_+)$, and bounded there by a constant depending
  on $\|h\|_\infty$ and the number of primes in
  $\operatorname{supp}(h)$;
\item symmetric (and in fact self-adjoint on its natural domain),
  by construction as a $\sigma$-symmetrized combination;
\item commutes with the $\sigma$-involution:
  $T_\zeta[h] \circ \sigma = \sigma \circ T_\zeta[h]$ by
  construction.
\end{enumerate}
We do \emph{not} claim that $T_\zeta[h]$ is finite-rank: it is a
finite sum of bounded dilation-plus-multiplication operators, each
of which has infinite rank on $L^2(\R_+, dx/x)$. Trace-class /
compact properties require an additional truncation (e.g.\ a cutoff
in $x$ compressing $L^2$ to a finite-dimensional subspace) which
is not supplied here; for the Hilbert--P\'olya programme, the
trace-class property is part of the open hypothesis $\mathbf{H_3}$.
\end{theorem}

\begin{proof}
(1) $C_c^\infty(\R_+)$ is dense in $L^2(\R_+, dx/x)$ by standard
analysis.

(2) Symmetry: for $f, g \in C_c^\infty(\R_+)$,
\[
  \langle V_{\sigma, h} f, g \rangle
  = \sum_p h(\log p) \frac{\log p}{p^{1/2}}
    \int_{\R_+} f(p x) \overline{g(x)} \frac{dx}{x}.
\]
Substituting $y = p x$, $dy/y = dx/x$:
\[
  = \sum_p h(\log p) \frac{\log p}{p^{1/2}}
    \int_{\R_+} f(y) \overline{g(y/p)} \frac{dy}{y}
  = \langle f, V_{\sigma, h}^* g \rangle,
\]
where $V_{\sigma, h}^* g(x) = \sum_p h(\log p) (\log p / p^{1/2})
g(x/p)$. Because $h$ is even, $h(\log p) = h(-\log p) = h(\log(1/p))$;
combined with the sum over $p$ (which is already self-adjoint under
the substitution $p \to 1/p$ when extended to the
$\sigma$-symmetric version), $V_{\sigma, h}^* = V_{\sigma, h}$
under a standard $\sigma$-symmetrization (explicit normalization in
Definition~\ref{def:Tzeta-h} assumed).

Symmetry of $T_0 \chi_h$: $T_0$ is self-adjoint on its natural core
(formal adjoint agrees by the same computation), and $\chi_h$ is a
multiplication operator by a real-valued bounded function, hence
self-adjoint.

(3) The $\sigma$-intertwining is the explicit functional-equation
symmetry: $\sigma \colon f(x) \mapsto f(1/x)$ maps the sum
$\sum_p h(\log p) (\log p / p^{1/2}) f(p x)$ to
$\sum_p h(\log p) (\log p / p^{1/2}) f(1/(p x)) = \sum_p h(\log p)
(\log p / p^{1/2}) f(x/p) = V_{\sigma, h} \sigma f$, since
$h$ is even.

Boundedness: since $h$ has compact support, only finitely many
primes $p$ satisfy $h(\log p) \neq 0$, so $V_{\sigma, h}$ is a
finite sum of bounded scaling operators on $L^2(\R_+, dx/x)$.
Each scaling $f \mapsto f(p x)$ is a bounded linear operator of
operator norm $1$ (change of variable is unitary-up-to-$p^{1/2}$
rescaling absorbed in the $p^{-1/2}$ factor); $T_0 \chi_h$ is
bounded on its natural domain. Hence $T_\zeta[h]$ is bounded, with
norm controlled by $\|h\|_\infty$ times the finite number of
primes in $\operatorname{supp}(h)$. \emph{We do not claim
finite-rank or trace-class}: a finite sum of nonzero dilation
operators is generically infinite-rank on $L^2(\R_+, dx/x)$.
\end{proof}

\begin{remark}[Intertwining and the critical line: a formal statement, not a spectral theorem]
\label{rmk:intertwining-critical}
The $\sigma$-intertwining property (3) of
Theorem~\ref{thm:tzeta-h} is the precise cascade-theoretic
realization of the $s \leftrightarrow 1-s$ functional-equation
symmetry of $\zeta$: under the Mellin transform, inversion
$f(x) \to f(1/x)$ corresponds to $s \to 1-s$. In the Hilbert--
P\'olya picture one would hope that a self-adjoint operator
intertwining with this symmetry has its Mellin-side
spectral-parameter representation reflection-symmetric about
$\Re s = 1/2$. This is a \emph{formal} structural statement about
the Mellin-side representation of the symmetry, \emph{not} a
proved concentration of eigenvalues on the critical line: the
latter would require control over the Mellin-transform
spectral calculus on $T_\zeta[h]$ not supplied here, and in any
case would not by itself prove the Riemann hypothesis (spectral
identification with the non-trivial zeros is part of the open
hypothesis $\mathbf{H_3}$).
\end{remark}

\subsection{The regularized Hecke transfer operator}
\label{subsec:TE-h}

Let $E$ be an elliptic curve over $\Q$ of conductor $N$, and let
$M_k(\Gamma_0(N), \Q) \otimes \C$ denote the space of weight-$k$
modular forms of level $\Gamma_0(N)$ with rational Fourier
coefficients; equip this with the Petersson inner product
\cite{DiamondShurman}.

\begin{definition}[Regularized Hecke transfer operator $T_E[h]$]
\label{def:TE-h}
For $h$ as in Definition~\ref{def:Vsigma-h}, define
\[
  T_E[h]
  := \sum_{p \nmid N} h(\log p) \cdot p^{-1/2} \cdot
    (T_p - a_p(E) \cdot \mathrm{id}),
\]
where $T_p$ is the classical Hecke operator on
$M_k(\Gamma_0(N))$ and $a_p(E) \in \Z$ is the $p$-th Fourier
coefficient of the modular form associated to $E$ by the
modularity theorem \cite{Wiles, BCDT}.
\end{definition}

\begin{theorem}[$T_E[h]$: bounded and well-defined]
\label{thm:tE-h}
For $h$ as above, $T_E[h]$ is a well-defined bounded endomorphism
of $M_k(\Gamma_0(N)) \otimes \C$. It is self-adjoint with respect
to the Petersson inner product and commutes with the other Hecke
operators $T_q$, $q \nmid N$.
\end{theorem}

\begin{proof}
Since $h$ has compact support, only finitely many primes
contribute, so $T_E[h]$ is a finite sum of bounded operators
$p^{-1/2}(T_p - a_p \cdot \mathrm{id})$. Each $T_p$ is
self-adjoint on the Petersson inner product (classical Hecke
theory \cite{DiamondShurman}), and the constant $a_p$ is real, so
$T_p - a_p \cdot \mathrm{id}$ is self-adjoint. Hecke operators at
different primes commute \cite{DiamondShurman}, giving commutation
with other $T_q$.
\end{proof}

\begin{remark}[Relation to the BSD $L$-function]
The sum
$\sum_p (\log p/p^{1/2})(T_p f - a_p f)$ at $f$ a Hecke eigenform
for $E$ would, if convergent, produce a scalar multiple of the
Dirichlet-series value $L'(E, s=1)/L(E, s=1)$. The cutoff $h$ keeps
the sum finite; removing the cutoff requires
$\mathbf{H_3}$ below.
\end{remark}

\subsection{Conditional trace matching under $\mathbf{H_3}$}
\label{subsec:H3}

\begin{hypothesis}[$\mathbf{H_3}$: Spectral--arithmetic trace
matching]
\label{hyp:H3}
There exists a sequence of cutoffs $h_N$ with
$h_N(\log p) \to 1$ pointwise as $N \to \infty$ and
$|h_N| \leq 1$, such that:
\begin{enumerate}
\item \textbf{Zeta case.} The renormalized traces
  $\operatorname{tr}_{\rm ren}(T_\zeta[h_N])$, defined by a
  regularization procedure subtracting the divergent
  $\sum_p \log p / p^{1/2}$ tail, converge as $N \to \infty$ to a
  limit matching the Weil explicit formula for $\zeta$-zeros
  \cite[\S 2.4]{DGOTI}.
\item \textbf{BSD case.} The renormalized traces
  $\operatorname{tr}_{\rm ren}(T_E[h_N])$ converge as
  $N \to \infty$ to a limit matching the explicit formula for
  $L(E, s)$-zeros
  \cite[\S 5]{IwaniecKowalski}.
\end{enumerate}
\end{hypothesis}

\begin{remark}[Status of $\mathbf{H_3}$]
$\mathbf{H_3}$ is \emph{open}. The explicit formulas for $\zeta$
and $L(E, s)$ \cite[\S 2.4]{DGOTI} give the form of the expected
limits; what $\mathbf{H_3}$ adds is that cascade cutoff removal
produces precisely those limits. Discharging $\mathbf{H_3}$ is the
core analytic content of RH and the analogous $L(E, 1)$-zero
conjectures.
\end{remark}

\begin{theorem}[$T_\zeta[h]$ and $T_E[h]$: conditional trace matching only]
\label{thm:trace-match-cond}
Assume $\mathbf{H_3}$. Then:
\begin{enumerate}
\item The renormalized \emph{traces}
  $\operatorname{tr}_{\rm ren}(T_\zeta[h_N])$ converge as
  $N \to \infty$ to a limit matching the Weil explicit formula
  for $\zeta$-zeros (by $\mathbf{H_3}$(i)). We do \emph{not}
  claim spectral (eigenvalue) convergence here, nor that all zeros
  lie on $\Re s = 1/2$: trace convergence is strictly weaker than
  spectral identification, and the $\sigma$-intertwining of
  Theorem~\ref{thm:tzeta-h}(3) gives a formal Mellin-side
  symmetry statement (Remark~\ref{rmk:intertwining-critical}),
  not a proved spectral consequence.
\item The renormalized traces $\operatorname{tr}_{\rm ren}(T_E[h_N])$
  converge to the $L(E, s)$-zero explicit-formula limit (by
  $\mathbf{H_3}$(ii)), with the same caveat that this is a trace
  statement, not a spectral / critical-line statement.
\end{enumerate}
\end{theorem}

\begin{proof}[Proof (sketch)]
Under $\mathbf{H_3}$, the trace convergence is by hypothesis. The
individual operators $T_\zeta[h_N], T_E[h_N]$ are self-adjoint
(Theorem~\ref{thm:tzeta-h}(2)); the $\sigma$-intertwining
(Theorem~\ref{thm:tzeta-h}(3)) combined with the Mellin-transform
identification of $\sigma$ with the functional-equation involution
gives a \emph{formal} reflection-symmetry of the Mellin-side
spectral-parameter representation about the relevant critical line
(see Remark~\ref{rmk:intertwining-critical}); we do not upgrade
this formal structural statement to a proved spectral theorem
about the operators' actual eigenvalues. We do
\emph{not} claim that this reflection-symmetry proves
concentration of the limit spectrum on $\Re s = 1/2$ or $\Re s =
1$: that would require, at minimum, spectral convergence (not just
trace convergence) plus simplicity arguments that $\mathbf{H_3}$
does not supply. A Clay-level proof would additionally require the
full regularization procedure and spectral convergence estimates
beyond $\mathbf{H_3}$, which are not discharged here.
\end{proof}

\begin{remark}[Scope]
Theorem~\ref{thm:trace-match-cond} does not resolve RH or BSD, and
does not establish a spectral-level correspondence to the zero
sets. It gives a \emph{trace-level} correspondence from
regularized cascade transfer operators to the explicit-formula
side of $\zeta$- and $L(E, s)$-zero distributions, under
$\mathbf{H_3}$. Any spectral statement (eigenvalue convergence,
critical-line concentration) requires strictly stronger hypotheses
than $\mathbf{H_3}$.
\end{remark}

%=====================================================================
\section{Cohomological and Computational Correspondences}
\label{sec:cohomology-comp}
%=====================================================================

This section constructs the Galois fixed-part projector
$\Pi_\sigma$ on $\Q(\pphi)$-valued cohomology (for Hodge) and the
transport between cascade Turing machines (CTMs) and ordinary
Turing machines (for $P$ vs $NP$). Unlike the previous three
sections, both constructions here are unconditional at the
theorem level --- the conditional parts concern universality
($\mathbf{H_4}$) rather than well-definedness.

\subsection{The Galois fixed-part projector $\Pi_\sigma$}
\label{subsec:pisigma}

\begin{definition}[Cohomology target]
\label{def:coh-target}
For a smooth projective complex variety $X$, let
$H^*(X, \Q(\pphi)) := H^*(X, \Q) \otimes_\Q \Q(\pphi)$ denote the
$\Q(\pphi)$-cohomology. The Galois action $\sigma$ acts coefficientwise:
$\sigma(\alpha \otimes \lambda) := \alpha \otimes \sigma(\lambda)$
for $\alpha \in H^*(X, \Q)$ and $\lambda \in \Q(\pphi)$.
\end{definition}

\begin{definition}[The projector $\Pi_\sigma$]
\label{def:pisigma}
Define
$\Pi_\sigma \colon H^*(X, \Q(\pphi)) \to H^*(X, \Q(\pphi))$ by
\[
  \Pi_\sigma(\omega) := \tfrac{1}{2} \big( \omega + \sigma(\omega)
  \big).
\]
This is a $\Q$-linear projection (since $\sigma^2 = \mathrm{id}$
and $\sigma$ is $\Q$-linear) whose image is the $\sigma$-fixed
subspace.
\end{definition}

\begin{theorem}[$\sigma$-fixed part equals rational part]
\label{thm:fix-sigma-rational}
For any smooth projective complex variety $X$,
\[
  \Fix(\sigma; H^*(X, \Q(\pphi))) = H^*(X, \Q) \otimes_\Q 1_{\Q(\pphi)}.
\]
In particular, the image of $\Pi_\sigma$ is (canonically isomorphic
to) $H^*(X, \Q)$, and $\Pi_\sigma$ is a $\Q$-linear projection onto
this subspace.
\end{theorem}

\begin{proof}
Write $\omega \in H^*(X, \Q(\pphi))$ as $\omega = \alpha + \beta
\pphi$ for unique $\alpha, \beta \in H^*(X, \Q)$ (using
$\Q(\pphi) = \Q \oplus \Q \pphi$ and the tensor decomposition).
Apply $\sigma$: $\sigma(\omega) = \alpha + \beta
\sigma(\pphi) = \alpha + \beta \bar\pphi = \alpha + \beta (1 -
\pphi) = (\alpha + \beta) - \beta \pphi$.

Equality $\omega = \sigma(\omega)$ gives: $\alpha + \beta \pphi =
(\alpha + \beta) - \beta \pphi$, i.e.\ $\beta = 0$ and $\alpha =
\alpha$. So $\Fix(\sigma) = \{\alpha + 0 \cdot \pphi : \alpha
\in H^*(X, \Q)\} = H^*(X, \Q)$, as claimed.
\end{proof}

\begin{theorem}[$\Pi_\sigma$ is functorial under cascade embeddings]
\label{thm:pisigma-functorial}
If $f \colon X \to Y$ is a morphism of smooth projective complex
varieties, then the pullback $f^* \colon H^*(Y, \Q(\pphi)) \to
H^*(X, \Q(\pphi))$ commutes with $\Pi_\sigma$:
$f^* \circ \Pi_\sigma = \Pi_\sigma \circ f^*$.
\end{theorem}

\begin{proof}
Standard functoriality: the pullback is $\Q$-linear and commutes
with the tensor-product structure on coefficients
\cite{GriffithsHarris}; in particular it commutes with the
coefficientwise Galois action $\sigma$. The projector $\Pi_\sigma
= (\mathrm{id} + \sigma)/2$ then commutes with $f^*$ by
$\Q$-linearity.
\end{proof}

\begin{remark}[On the Hodge Millennium application]
Theorem~\ref{thm:fix-sigma-rational} gives the precise version of
the $\sigma$-fixed / rational identification used in
\cite{MillenniumHodge}. The upgrade to the Hodge Conjecture (every
rational $(p, p)$ class is algebraic) requires additional
universality / embeddability machinery, which is the content of
$\mathbf{H_{4a}}$ below.
\end{remark}

\subsection{Cascade Turing machines and transport}
\label{subsec:ctm-tm}

\begin{definition}[Cascade Turing machine (CTM)]
\label{def:ctm}
A \emph{cascade Turing machine} is a tuple $(Q, \Sigma, \Gamma,
\delta_{\rm cas}, q_0, q_{\rm acc}, q_{\rm rej})$, where:
\begin{itemize}
\item $Q$ is a finite set of states,
\item $\Sigma \subseteq \Gamma$ is the input alphabet, with
  $\Gamma$ the tape alphabet;
\item $q_0, q_{\rm acc}, q_{\rm rej} \in Q$ are start, accept, and
  reject states;
\item $\delta_{\rm cas} \colon Q \times \Gamma \to Q \times \Gamma
  \times \{L, R, {\rm F}, \sigma, {\rm R}, P_r\}$ is a transition
  function that specifies, in addition to standard TM moves
  $\{L, R\}$ (left/right), the four cascade primitives
  $\{\rm F, \sigma, R, P_r\}$ defined in \cite[\S D1]{CascadeDynamics}:
  $\rm F$ (closure evaluation), $\sigma$ (Galois involution),
  $\rm R$ (rung restriction), $P_r$ (refinement step).
\end{itemize}
A CTM computation is defined as in the classical TM model, with
each $\rm F, \sigma, R, P_r$ step interpreted as a specific
finite-time cascade operation on a cascade-data auxiliary
register.
\end{definition}

\begin{definition}[Transport maps]
\label{def:ctm-transport}
Define:
\begin{itemize}
\item the \emph{encoding} $\Pi_{\TM} \colon {\rm CTM} \to {\rm TM}$
  sending a CTM $M$ to a classical TM
  $\Pi_{\TM}(M)$ that simulates $M$ by expanding each cascade
  primitive into a fixed number of classical TM steps.
  The expansion is defined per primitive: $\rm F$ becomes a
  polynomial-time classical computation of closure energy on the
  finite-level discrete substrate (Section~\ref{sec:dynamics});
  $\sigma$ becomes a polynomial-time classical $\Q(\pphi)$-arithmetic
  operation; $\rm R$ and $P_r$ become finite-time combinatorial
  operations.
\item the \emph{decoding} $I_{\TM} \colon {\rm TM} \to {\rm CTM}$
  sending a classical TM $T$ to a CTM
  $I_{\TM}(T)$ that uses none of the cascade primitives (so
  $\delta_{\rm cas}$ restricts to the classical
  $\delta \colon Q \times \Gamma \to Q \times \Gamma \times
  \{L, R\}$).
\end{itemize}
\end{definition}

\begin{theorem}[CTM/TM transport: decidability; polynomial-time preservation only under uniform-cost hypothesis]
\label{thm:ctm-tm-transport}
\begin{enumerate}
\item (Decidability.) A language $L$ is decidable by a CTM if and
  only if $L$ is decidable by a classical TM. In particular,
  $\Pi_{\TM}$ and $I_{\TM}$ preserve the decidability class.
\item (Polynomial-time, \emph{conditional on a uniform-cost
  hypothesis}.) Assume: there exist a polynomial $q$ and a
  constant $k$ such that, for every CTM configuration $C$ of size
  $\leq n$, the classical-TM expansion of each cascade primitive
  applied at $C$ runs in time $\leq q(n)^k$. Under this hypothesis,
  if a CTM $M$ decides $L$ in cascade time $T_M(n)$, then
  $\Pi_{\TM}(M)$ decides $L$ in classical time $T_M(n) \cdot
  q(T_M(n))^k$, polynomial in $T_M(n)$; in particular $L \in
  \mathrm{P}_{\rm CTM}$ iff $L \in \mathrm{P}_{\rm TM}$. Without
  the uniform-cost hypothesis, primitive-expansion costs depend on
  the current configuration size (e.g.\ discrete-Laplacian cost
  grows with substrate size), and a size-independent uniform
  constant $c$ cannot in general be extracted.
\end{enumerate}
\end{theorem}

\begin{proof}
(1) Forward direction: given a CTM $M$, the classical TM
$\Pi_{\TM}(M)$ simulates $M$ by expanding each primitive as
specified in Definition~\ref{def:ctm-transport}. Each primitive is
a finite classical computation, so the simulation halts exactly
when $M$ halts, on the same input, to the same verdict.

Reverse: given a classical TM $T$, $I_{\TM}(T)$ is a CTM that uses
only $\{L, R\}$ primitives, simulating $T$ identically.

(2) Under the uniform-cost hypothesis: each CTM step at
configuration size $\leq n$ is simulated in $\leq q(n)^k$
classical-TM steps, so a CTM computation of $T_M(n)$ steps is
simulated in $\leq T_M(n) \cdot q(T_M(n))^k$ classical-TM steps,
polynomial in $T_M(n)$ if $q, k$ are fixed. Without the
uniform-cost hypothesis, the simulation cost of the $t$-th
primitive depends on the configuration size $s_t$ at that point
(e.g.\ $\rm F$ computes a discrete Laplacian on a finite substrate
whose size grows with $t$), and $\sum_{t \leq T_M(n)} q(s_t)^k$
need not admit a polynomial bound in $T_M(n)$ alone.
\end{proof}

\subsection{Conditional universality under $\mathbf{H_4}$}
\label{subsec:H4}

\begin{hypothesis}[$\mathbf{H_{4a}}$: Cascade embeddability
(universality for Hodge)]
\label{hyp:H4a}
Every smooth projective complex variety $X$ admits a
\emph{cascade-compatible embedding} $\iota_X$ defined in
\cite[\S 4--6]{CascadeEmbeddings} (formal-framework grade in the
cited source), such that the Hodge class structure on
$H^*(X, \Q)$ is compatible with the $\sigma$-fixed-part projector
$\Pi_\sigma$ on $H^*(X, \Q(\pphi))$ under pullback by $\iota_X$.
\end{hypothesis}

\begin{hypothesis}[$\mathbf{H_{4b}}$: CTM normal form (universality
for $P$ vs $NP$)]
\label{hyp:H4b}
Every polynomial-time CTM algorithm admits a \emph{cascade normal
form} in which the cascade primitives $\{\rm F, \sigma, R, P_r\}$
are used in a restricted sequence, allowing the $\sigma$-irreducibility
property of \cite[\S D5--D10]{CascadeDynamics} (formal-framework
grade) to distinguish $\mathrm{P}_{\rm CTM}$ from $\mathrm{NP}_{\rm CTM}$.
\end{hypothesis}

\begin{remark}[Status of $\mathbf{H_4}$]
Both $\mathbf{H_{4a}}$ and $\mathbf{H_{4b}}$ are \emph{open}.
$\mathbf{H_{4a}}$ is the Hodge-universality claim that the
cascade-embeddability of \cite[\S 6]{CascadeEmbeddings} extends from
the specific examples there (complex projective spaces, flag
varieties, abelian varieties with CM, certain complete
intersections) to all smooth projective complex varieties; this is
listed in \cite[\S MP5]{CascadeMillennium} as an open technical
item. $\mathbf{H_{4b}}$ is the cascade-algorithm normal-form
hypothesis used in the earlier \cite{MillenniumPNP} (then called
``cascade structural completeness''); it is a cascade-internal
combinatorial conjecture with no classical analogue yet.
\end{remark}

\begin{theorem}[Hodge and $P$ vs $NP$ transport, conditional on
$\mathbf{H_4}$]
\label{thm:universality-cond}
\begin{enumerate}
\item Under $\mathbf{H_{4a}}$: for every smooth projective complex
  variety $X$, the Hodge classes $H^{p, p}(X, \Q) \subset H^*(X, \Q)$
  are canonically isomorphic to $\Fix(\sigma;
  H^*(X, \Q(\pphi)))_{(p, p)}$, the degree-$(p,p)$ part of the
  $\sigma$-fixed cohomology. In particular, every $\sigma$-fixed
  rational $(p, p)$ class pulls back under $\iota_X^*$ to a class
  on $X$ in the image of $\Pi_\sigma$.
\item Under $\mathbf{H_{4b}}$: the class $\mathrm{P}_{\rm CTM}^\sigma$
  of $\sigma$-preserving polynomial-time CTM algorithms is a
  proper subset of $\mathrm{NP}_{\rm CTM}$, and under the
  CTM-TM transport of Theorem~\ref{thm:ctm-tm-transport}(2), this
  descends to $\mathrm{P}_{\rm TM} \neq \mathrm{NP}_{\rm TM}$.
\end{enumerate}
\end{theorem}

\begin{proof}
Both statements are immediate from combining their respective
hypotheses with the unconditional theorems of
Sections~\ref{subsec:pisigma} (for $\Pi_\sigma$) and
\ref{subsec:ctm-tm} (for $\Pi_{\TM}, I_{\TM}$). The role of
$\mathbf{H_4}$ is to provide the universal extension to arbitrary
varieties (respectively arbitrary polynomial-time algorithms) that
the unconditional theorems do not cover.
\end{proof}

\begin{remark}[Scope]
Theorem~\ref{thm:universality-cond} does not resolve Hodge or $P$
vs $NP$. It gives the structural transport from cascade-internal
correspondences to classical statements, under universality /
normal-form hypotheses that are themselves open.
\end{remark}

%=====================================================================
\section{Compatibility Lemmas}
\label{sec:compat}
%=====================================================================

This short section records the commutative-diagram identities
tying the three commuting symmetry actions (Lemma~\ref{lem:commute})
to the seven correspondences of
Sections~\ref{sec:metric-hydro}--\ref{sec:cohomology-comp}.

\begin{proposition}[Projection/refinement compatibility]
\label{prop:proj-refine-compat}
For each
$\Pi \in \{\Pi_{\met}, \Pi_{\hyd}, \Pi_{\gge}, T_\zeta[h], T_E[h],
\Pi_\sigma, \Pi_{\TM}\}$ and every $n \geq 0$, the refinement
shift $R_n$ commutes with $\Pi^{(n)}$ up to the target-inclusion
maps introduced in Theorem~\ref{thm:abstract-correspondence}(A4).
\end{proposition}

\begin{proof}
$\Pi_{\met}, \Pi_{\hyd}, \Pi_{\gge}$ commute with refinement by the
zero-extension property: moment tensors, averaged fields, and
holonomies on zero-extended configurations vanish on the new edges
and vertices. $T_\zeta[h], T_E[h], \Pi_\sigma, \Pi_{\TM}$ are
level-independent, acting on $L^2(\R_+, dx/x)$,
$M_k(\Gamma_0(N))$, $H^*(X, \Q(\pphi))$, and
$\mathrm{CTM} \leftrightarrow \mathrm{TM}$ respectively, none of
which is indexed by the refinement level.
\end{proof}

\begin{proposition}[Projection/$\sigma$ compatibility]
\label{prop:proj-sigma-compat}
$\Pi_{\met}, \Pi_{\hyd}, \Pi_{\gge}$ are $\sigma$-invariant.
$\Pi_\sigma$ is $\sigma$-equivariant with trivial quotient.
$T_\zeta[h]$ is $\sigma$-intertwined (Theorem~\ref{thm:tzeta-h}(3)),
and $T_E[h]$ is intertwined with the Atkin--Lehner involution
\cite[\S 5.5]{DiamondShurman}.
$\Pi_{\TM}$ is $\sigma$-invariant since its target is classical
TMs on which $\sigma$ acts trivially.
\end{proposition}

\begin{proof}
Each assertion is a direct verification using the sector-structure
of $\sigma$ (Definition~\ref{def:sigma-action-subst}) or restates a
theorem proved earlier.
\end{proof}

\begin{proposition}[Projection/dynamics compatibility]
\label{prop:proj-dynamics-compat}
The finite-level cascade gradient flow $e^{-t L_n}$ commutes with
each projection up to conditional limits:
\begin{enumerate}
\item Under $\mathbf{H_1}$ \emph{together with} the P5--P7
  structural imports (Hypotheses~\ref{hyp:coord-inj-import},
  \ref{hyp:tiling-import}, \ref{hyp:separation-import},
  \ref{hyp:gradient-image-import}) and, for the full-flow case,
  $\mathbf{H_1'}$ (Hypothesis~\ref{hyp:full-flow}): the flow
  commutes with $\Pi_{\met}, \Pi_{\hyd}$, with target operators
  the Ricci-type operator and the NS operator respectively
  (Theorems~\ref{thm:pimet-cond}, \ref{thm:pihyd-cond}). Without
  $\mathbf{H_1'}$, only the mass-only sub-case of the flow
  commutation survives.
\item Under $\mathbf{H_2}$ \emph{together with} the modeling
  choices (M1)--(M3) of \S\ref{subsec:gauge-modeling}: the flow
  commutes with $\Pi_{\gge}$, with target the YM heat flow
  (Theorem~\ref{thm:pigauge-cond}).
\item Unconditionally: on the $H_4$ sector the flow induces a
  bounded-operator perturbation of $T_\zeta[h], T_E[h]$ preserving
  self-adjointness.
\end{enumerate}
\end{proposition}

\begin{proof}
(1)--(2) are contained in the cited theorems.
(3) follows from bounded-operator continuity of $e^{-t L_n}$ on
$L^2(\R_+, dx/x)$ restricted to the compactly-supported spectral
sector.
\end{proof}

%=====================================================================
\section{Worked Examples and Sanity Checks}
\label{sec:examples}
%=====================================================================

\subsection{$24$-cell point-source moment tensor}

Proposition~\ref{prop:moment-point-source} computes
$M_0(\Phi^{\rm pt})(v) \propto \hat r(v) \otimes \hat r(v)$; this
matches \cite[\S C4.1]{CascadeGR}.

\subsection{Octonionic Wilson holonomy on a $24$-cell face}

\begin{example}
\label{ex:octo-holonomy}
Let $A \in X_0^1$ satisfy $A(e_i) = t_i i_0$ on the four oriented
edges of a square face, zero elsewhere. Then
$\Pi_{\gge}^{(0)}(A)(e_i) = \exp(a_0 t_i \cdot
P_{\mathfrak{su}(3)}(\mathrm{ad}_{i_0}))$. The loop holonomy is, to
leading order in $a_0$, $\exp(a_0(t_1+t_2+t_3+t_4)
P_{\mathfrak{su}(3)}(\mathrm{ad}_{i_0}))$, the standard Wilson-loop
form \cite{WilsonLatticeGauge}.
\end{example}

\subsection{$\sigma$-fixed cohomology of $\CP^1$}

\begin{example}
\label{ex:cp1-sigma}
$H^*(\CP^1, \Q) = \Q[0]\oplus\Q[2]$, $\Q$-dimension $2$.
Tensoring with $\Q(\pphi)$ gives $H^*(\CP^1, \Q(\pphi))$ of
$\Q$-dimension $4$, with $\sigma$-fixed part
$\Q[0]\oplus\Q[2]$ of $\Q$-dimension $2$. This confirms
Theorem~\ref{thm:fix-sigma-rational} for $\CP^1$.
\end{example}

\subsection{CTM encoding of binary addition}

\begin{example}
\label{ex:ctm-add}
Any classical $O(n)$-step TM for binary addition lifts via
$I_{\TM}$ to a CTM of the same size and time, and
$\Pi_{\TM} \circ I_{\TM}$ is the identity on classical TMs. This
confirms that the transport is the identity on the sub-class of
TMs, as the CTM model extends the classical TM model.
\end{example}

%=====================================================================
\section{Relation to the Six Millennium Papers}
\label{sec:handoff}
%=====================================================================

\begin{table}[h]
\caption{Handoff to the six Clay Millennium companion papers, plus
Poincar\'e.}
\label{tab:handoff}
\begin{tabular}{p{1.15in}p{1.8in}p{1.3in}p{1.4in}}
\hline
Paper & Cites from this paper & Gap closed & What remains open \\
\hline
\cite{MillenniumYM} &
  Thm.~\ref{thm:pigauge-n}; Thm.~\ref{thm:pigauge-cond} &
  precise $\Pi_{\gge}$; gauge action transport &
  $\mathbf{H_2}$: OS axioms, area law \\
\cite{MillenniumRH} &
  Def.~\ref{def:Tzeta-h}; Thms.~\ref{thm:tzeta-h},
  \ref{thm:trace-match-cond}(1) &
  reg.\ $T_\zeta[h]$; $\sigma$-intertwined &
  $\mathbf{H_3}$: cutoff-removal match \\
\cite{MillenniumNS} &
  Def.~\ref{def:pihyd-n}; Thms.~\ref{thm:pihyd-n},
  \ref{thm:pihyd-cond} &
  precise $\Pi_{\hyd}$; coarse-graining &
  $\mathbf{H_1}$: BI+CC+overshoot \\
\cite{MillenniumHodge} &
  Def.~\ref{def:pisigma}; Thms.~\ref{thm:fix-sigma-rational},
  \ref{thm:pisigma-functorial} &
  $\sigma$-fixed=rational; functoriality &
  $\mathbf{H_{4a}}$: cascade-embeddability \\
\cite{MillenniumBSD} &
  Def.~\ref{def:TE-h}; Thms.~\ref{thm:tE-h},
  \ref{thm:trace-match-cond}(2) &
  reg.\ $T_E[h]$; Hecke-compat.\ &
  $\mathbf{H_3}$: $L(E,s)$-trace match \\
\cite{MillenniumPNP} &
  Def.~\ref{def:ctm}; Thms.~\ref{thm:ctm-tm-transport},
  \ref{thm:universality-cond}(2) &
  precise CTM/TM transport &
  $\mathbf{H_{4b}}$: cascade normal form \\
\cite{MillenniumPoincare} &
  Thm.~\ref{thm:pimet-cond} &
  precise continuum metric proj.\ &
  none additional (Perelman closed) \\
\hline
\end{tabular}
\end{table}

\begin{remark}[Reading the table]
Each Millennium paper inherits a hypothesis stack whose
\emph{headline} family is one of $\mathbf{H_1}$--$\mathbf{H_4}$.
The actual stack is larger: the metric / hydrodynamic entries
additionally inherit the P5--P7 structural imports
(Hypotheses~\ref{hyp:coord-inj-import}--\ref{hyp:gradient-image-import})
and, for the full-flow PDE portion, the foundations-level
conjecture $\mathbf{H_1'}$ (Hypothesis~\ref{hyp:full-flow}); the
gauge entry additionally inherits the modeling choices (M1)--(M3)
of \S\ref{subsec:gauge-modeling}. After citation, the open work
in each Millennium paper is the specific classical analytic or
algorithmic question captured by its hypothesis stack, not a
cascade-internal ambiguity about what the projection \emph{is}.
\end{remark}

\begin{remark}[What this paper does not hand off]
Closure-coefficient matching (\cite[\S F8]{CascadeFoundations}),
universality claims not captured by $\mathbf{H_4}$, and the
P4--P8 bootstrap identities of \cite{CascadePregeometry} are not
part of this paper's handoff; see Table~\ref{tab:exclusions}.
\end{remark}

\section{Conclusion}
\label{sec:conclusion}

The systemic gap in the earlier Millennium drafts was the lack of
precise mathematical definitions for the cascade-to-classical
projections. This paper closes that gap by constructing seven
explicit maps ($\Pi_{\met}, \Pi_{\hyd}, \Pi_{\gge}, T_\zeta[h],
T_E[h], \Pi_\sigma, (\Pi_{\TM}, I_{\TM})$) with specified domain,
codomain, continuity, and equivariance. Finite-level and
finite-cutoff versions are proved unconditionally well-defined.
Three continuum-limit extensions (metric, hydrodynamic, gauge) and
the spectral-arithmetic trace matching are conditional theorems
under four open hypothesis families $\mathbf{H_1}$--$\mathbf{H_4}$.
The cohomological and CTM/TM correspondences are proved
unconditionally at the level the companion papers actually need.

The paper does not prove any Millennium Prize Problem, does not
resolve any cascade-internal technical item
($\mathbf{H_1}$--$\mathbf{H_4}$ remain open), and explicitly
excludes the P4--P8 bootstrap and F8 closure-coefficient-matching
machinery. Its contribution is structural: replacing informal
``projection'' language with precise correspondences so the six
Millennium companion manuscripts can be evaluated at
Clay-submission rigour against a stable foundation.

\appendix
\section{Citation and Notation Audit}
\label{app:audit}

\subsection{Notation concordance}

\begin{table}[h]
\caption{Notation concordance.}
\label{tab:notation}
\begin{tabular}{lll}
\hline
This paper & Cascade-derivation sources & Millennium manuscripts \\
\hline
$\pphi$ & $\varphi$ \cite{CascadeFoundations} & $\varphi$ \\
$\sigma$ & $\sigma$ \cite[\S F5]{CascadeFoundations} & $\sigma$ \\
$V_{H_4}$ & $H_4$-rung fibre \cite[\S 3]{CascadeEmbeddings} & $H_4$ sector \\
$V_{D_4}$ & $D_4$-rung fibre \cite[\S 2]{CascadeEmbeddings} & $D_4$ sector \\
$P_{H_4}, P_{D_4}$ & $\pi_H$ (related to $\iota_D$) & $\pi_H, \iota_D$ \\
$X_n^0, X_n^1$ & vertex- / edge-sector at level $n$ (new) & implicit \\
$\mathcal{F}$ & $\alpha R + \beta E - \gamma Q$ \cite[\S F2]{CascadeFoundations} & closure functional \\
$F_n$ & finite-level closure energy (new) & --- \\
$\Pi_{\met}$ & informal ``metric projection'' \cite{MillenniumNS, MillenniumPoincare} & previously undefined \\
$\Pi_{\hyd}$ & informal ``coarse-graining'' \cite{MillenniumNS} & --- \\
$\Pi_{\gge}$ & informal ``gauge projection'' \cite{MillenniumYM} & --- \\
$T_\zeta[h]$ & reg.\ of ``$T_\zeta$'' \cite{MillenniumRH} & unreg.\ version excluded \\
$T_E[h]$ & reg.\ of ``$T_E$'' \cite{MillenniumBSD} & unreg.\ version excluded \\
$\Pi_\sigma$ & ``$\sigma$-fixed projection'' \cite{MillenniumHodge} & $\Pi_\sigma$ \\
$\Pi_{\TM}, I_{\TM}$ & CTM/TM equiv.\ \cite{MillenniumPNP} & --- \\
$\mathbf{H_i}$ & ``open items'' & conditional hypotheses \\
\hline
\end{tabular}
\end{table}

\subsection{Citation audit}

\begin{table}[h]
\caption{Citation audit (cascade-derivation sources).}
\label{tab:audit-citations}
\begin{tabular}{lll}
\hline
Citation & Used for & Content at cited location \\
\hline
\cite{CascadeFoundations} \S F1 & $\pphi$-uniqueness & proved \\
\cite{CascadeFoundations} \S F2 & closure form & proved \\
\cite{CascadeFoundations} \S F3 & 7-rung cascade & proved \\
\cite{CascadeFoundations} \S F5 & $\sigma$-involution & proved \\
\cite{CascadeEmbeddings} \S\S 2--3 & $D_4, H_4$ & explicit coords \\
\cite{CascadeEmbeddings} \S 5 & $\mathrm{Cl}(1,3)$ at rung 16 & explicit iso \\
\cite{CascadeEmbeddings} \S 6 & oct.\ at rung 8, $\mathrm{SU}(3) \subset G_2$ & explicit $V_8 \cong \OO$ \\
\cite{CascadeGR} \S C4.1 & moment formula & pt-source comp.\ \\
\cite{CascadeGR} \S C2.bis & GH continuum limit & statement + sketch (F) \\
\cite{CascadeGR} \S C3.bis & $A_5$-density & proof + transcendence (F) \\
\cite{CascadeDynamics} \S\S D1--D4 & CTM primitives & defs \\
\cite{CascadeDynamics} \S\S D5--D10 & complexity classes & formal framework \\
\cite{CascadeMillennium} \S MP6.3 & Poincar\'e status & ``no contribution beyond consistency'' \\
\hline
\end{tabular}
\end{table}

\begin{thebibliography}{99}

\bibitem{CascadeFoundations}
\emph{Cascade Foundations: F1--F8}, internal manuscript,
  \texttt{papers/cascade-derivation/cascade-foundations.md}.

\bibitem{CascadeEmbeddings}
\emph{Cascade Embeddings: $D_4 \subset E_8$, $H_4$ in $E_8$, rung
  fibres}, internal manuscript,
  \texttt{papers/cascade-derivation/cascade-embeddings.md}.

\bibitem{CascadeGR}
\emph{Cascade GR: Sub-phase Programme}, internal manuscript,
  \texttt{papers/cascade-derivation/cascade-gr.md}; moment tensor
  \S C4.1; GH limit (formal framework) \S C2.bis; $A_5$-density
  (formal framework) \S C3.bis.

\bibitem{CascadeDynamics}
\emph{Cascade Dynamics: Primitives and Complexity Classes},
  internal manuscript,
  \texttt{papers/cascade-derivation/cascade-dynamics.md}; CTM
  primitives \S\S D1--D4; complexity classes (formal framework)
  \S\S D5--D10.

\bibitem{CascadeMillennium}
\emph{Cascade and the Millennium Problems}, internal manuscript,
  \texttt{papers/cascade-derivation/cascade-millennium.md}.

\bibitem{CascadePregeometry}
\emph{Cascade Pregeometry: P0--P9} (formal framework),
  \texttt{papers/cascade-derivation/cascade-pregeometry.md}.

\bibitem{CascadeBridges}
\emph{Cascade Bridges},
  \texttt{papers/cascade-derivation/cascade-bridges.md}.

\bibitem{MillenniumYM}
\emph{Yang--Mills Mass Gap via Cascade Substrate},
  \texttt{papers/millennium-ym/ym-mass-gap.tex}.

\bibitem{MillenniumRH}
\emph{Riemann Hypothesis via Cascade Spectral Zeta},
  \texttt{papers/millennium-rh-formal/rh-formal.tex}.

\bibitem{MillenniumNS}
\emph{Navier--Stokes Global Smoothness via Cascade Coarse-Graining},
  \texttt{papers/millennium-ns-formal/ns-formal.tex}.

\bibitem{MillenniumHodge}
\emph{Hodge Conjecture via Cascade $\sigma$-Rationality},
  \texttt{papers/millennium-hodge-formal/hodge-formal.tex}.

\bibitem{MillenniumBSD}
\emph{BSD Conjecture via Cascade Hecke Operator},
  \texttt{papers/millennium-bsd-formal/bsd-formal.tex}.

\bibitem{MillenniumPNP}
\emph{P vs NP via Cascade Dynamics},
  \texttt{papers/millennium-pnp-formal/pnp-formal.tex}.

\bibitem{MillenniumPoincare}
\emph{Cascade Compatibility with the Poincar\'e Theorem},
  \texttt{papers/millennium-poincare-formal/poincare-formal.tex}.

\bibitem{CodexReview}
\emph{Codex review log for Millennium manuscripts}, internal
  archive (iterations 1--3 of the Poincar\'e paper demonstrate
  the projection-rigor gap this paper addresses).

\bibitem{CoxeterRegularPolytopes}
H.~S.~M.~Coxeter, \emph{Regular Polytopes}, 3rd ed., Dover, 1973.

\bibitem{BaezOctonions}
J.~Baez, \emph{The Octonions}, Bull.\ AMS \textbf{39} (2002),
  145--205.

\bibitem{BuragoBuragoIvanov}
D.~Burago, Y.~Burago and S.~Ivanov, \emph{A Course in Metric
  Geometry}, Graduate Studies in Mathematics \textbf{33}, AMS, 2001.

\bibitem{CheegerColding1997}
J.~Cheeger and T.~H.~Colding, \emph{On the structure of spaces
  with Ricci curvature bounded below. I}, J.\ Diff.\ Geom.\
  \textbf{46} (1997), 406--480.

\bibitem{GriffithsHarris}
P.~Griffiths and J.~Harris, \emph{Principles of Algebraic
  Geometry}, Wiley, 1978.

\bibitem{WilsonLatticeGauge}
K.~G.~Wilson, \emph{Confinement of quarks}, Phys.\ Rev.\ D
  \textbf{10} (1974), 2445--2459.

\bibitem{OSBook}
K.~Osterwalder and R.~Schrader, \emph{Axioms for Euclidean
  Green's functions}, Comm.\ Math.\ Phys.\ \textbf{31} (1973),
  83--112.

\bibitem{DGOTI}
H.~Iwaniec, \emph{Lectures on the Riemann Zeta Function}, AMS
  University Lecture Series \textbf{62}, 2014.

\bibitem{IwaniecKowalski}
H.~Iwaniec and E.~Kowalski, \emph{Analytic Number Theory}, AMS
  Colloquium Publications \textbf{53}, 2004.

\bibitem{DiamondShurman}
F.~Diamond and J.~Shurman, \emph{A First Course in Modular Forms},
  Graduate Texts in Mathematics \textbf{228}, Springer, 2005.

\bibitem{Wiles}
A.~Wiles, \emph{Modular elliptic curves and Fermat's Last Theorem},
  Annals of Mathematics \textbf{141} (1995), 443--551.

\bibitem{BCDT}
C.~Breuil, B.~Conrad, F.~Diamond, and R.~Taylor, \emph{On the
  modularity of elliptic curves over $\Q$: wild 3-adic exercises},
  J.\ AMS \textbf{14} (2001), 843--939.

\bibitem{Bogachev}
V.~I.~Bogachev, \emph{Measure Theory}, Vol.\ I, Springer, 2007.

\bibitem{ReedSimonI}
M.~Reed and B.~Simon, \emph{Methods of Modern Mathematical
  Physics, I: Functional Analysis}, Academic Press, 1980.

\end{thebibliography}

\end{document}
